\begin{chapterbox}
    \chapter{Quruns Erwachen (2021)}
    \label{Quruns Erwachen (2021)}
    \az{Jahr 71}
    
    Der Seher Leander wird von der Stammesältesten der Agren ins Graue Gebirge berufen. Ihr mysteriöser Auftrag schickt gleich mehrere mächtige Entitäten auf Kollisionskurs. Sie alle sind auf der Suche nach etwas. Nach einem alten Kameraden. Nach Freiheit. Nach Seelennahrung. Nach einem Platz in dieser Welt. Nach einer Möglichkeit, einen Rivalen zu beseitigen. Und... nach dem perfekten Kessel aus Zwergenstahl?
\end{chapterbox}


\section{Prolog}

\az{Jahr 71}

Der Raum war kreisrund und leer bis auf einen Altar in seiner Mitte und eine riesige Gestalt am gegenüberliegenden Ende des Raums. Dort lag ein versteinerter Drache, die Überreste des einst so edlen Kardòl, den von Rissen durchzogenen Kopf in seinen stacheligen Schwanz verbissen. Aber ich war nicht seinetwegen gekommen.

Rotes Mondlicht fiel vom Himmel durch einen kleinen Riss in der Höhlendecke auf den kalten Fels. Es beleuchtete den Altar ebenso wie den weißen Edelstein, den ich soeben darauf gelegt hatte. Ich konnte den Raum und die Farben natürlich nicht sehen, aber ich hatte in den letzten Tagen genug oft davon geträumt, um sie mir gut vorstellen zu können.

Dieser weiße Edelstein war außergewöhnlich. „Drachenherz“ hatten die Zwerge ihn genannt und ihn aufgrund seiner Größe für etwas derart Besonderes gehalten, dass sie mich zweifelsohne für den Rest meines Lebens in den Tiefminen schuften lassen würden, sollten sie je herausfinden, dass ich ihn mir angeeignet hatte. Narren waren sie! Versehen mit den richtigen Zwergenrunen, besaß der Kristall vor mir nun viel interessantere Eigenschaften als seine Reinheit und Größe. Eigenschaften magischer Natur. Selbst jemand wie ich, der noch nie ein Talent für Zaubersprüche und mächtigeren Hokuspokus gezeigt hatte, konnte diese Eigenschaften zu seinem Vorteil nutzen.

Ich fühlte die Präsenz des Bösen, noch ehe es seinen Mund öffnete. Meine Nackenhaare stellten sich auf und ein Schauer lief über meinen gesamten Körper. Ich versuchte, mir nichts anmerken zu lassen, rückte meine Augenbinde zurecht und schritt langsam auf den Altar zu.

„Du hast mich gerufen?“, sprach eine tiefe, heisere Stimme aus den Tiefen der Erde. Wenn ich mich nicht irrte, waren das Worte aus der dunklen Sprache Krahds. Es gibt keine Sprache auf dieser Welt, derer ich nicht zumindest teilweise mächtig wäre. So stellte es für mich kein Problem dar, das Böse zu verstehen. Dennoch antwortete ich nicht und schaute verwirrt drein. Zu zeigen, dass ich die Sprache der Krahder beherrschte, würde nur unnötig Alarmglocken läuten.

„Wer hat mich gerufen?“, wiederholte die tiefe Stimme nun in der Sprache der Andori, mit leicht gebrochenem Akzent. Es schien, als würde der Urheber der Stimme die Sprache zwar kennen, sie aber zutiefst verabscheuen. Und das schien natürlich so, weil es so war.

Andächtig antwortete ich: „Es ist also wahr! Nicht alle Geister der Gefallenen ziehen weiter. Ihr seid der lebendige Beweis für meine These! Nun, so lebendig, wie es den Umständen entsprechend geht. Gepriesen sei Mutter Natur!“

Die tiefe Stimme gluckste auf und ich grinste innerlich. Das Böse schien mir die Rolle des begeisterten Gelehrten abzukaufen.

„Hoffentlich bin ich nicht umsonst in diese Kammer tief ins Graue Gebirge gestiegen. Was willst du von mir, du blauer Wicht?“

Äußerlich ließ ich mich verunsichern, innerlich triumphierte ich. Stockend ließ ich die Worte meinen Mund verlassen: „Ich... ich dachte... nun... wenn ihr es geschafft habt, Gevatter Tod ein Schnippchen zu schlagen... vielleicht... vielleicht könntet ihr mich lehren? Es mir beibringen?“

Ich verstummte. Die Kunst, den Tod zu überdauern, ist eine äußerst nützliche Fähigkeit, die ich nur allzu gerne selbst besäße. Ich hatte meine Lebensdauer bereits mit Tränkchen und Süppchen aus der gesamten bekannten Welt verlängert, dennoch nahmen mit jedem Jahr die Falten in meinem Gesicht zu. Noch. Aber ich war nicht in Eile. Ich hatte Zeit. Sobald das Böse erst einmal in meiner Gewalt wäre, würde ich sein Wissen langsam, aber sicher aus ihm herauspressen können.

Das Böse lachte krächzend auf: „Hast du das Zeug, den Tod zu überlisten? Bist du bereit, all das, was du besitzt, dafür aufzugeben? Bist du bereit, jede mögliche Grenze zu überschreiten, die man überschreiten könnte?“

„Ja!“, antwortete ich mit gespielt zitternder Stimme, „Jede Grenze außer der zum Reich des Todes!“

„Guuuut“, zischte die Stimme des Bösen. „Dann habe ich eine Aufgabe für dich. Tritt ins Mondlicht, mein neuer Freund.“

Die Luft wurde kaum merklich wärmer, als ich ins Licht des roten Mondes trat und auf den Blutsteinaltar kletterte. Eine Spannung lag in der Luft und der Geruch nach Schweiß und Blut drang in meine Nase.

„Was nun?“

„Es ist ganz einfach. Alles, was du tun musst... ist sterben!“

Ein Windstoß fuhr durch die Höhle und wehte mich beinahe vom Altar. Natürlich hatte das Böse nicht vorgehabt, mir das Geheimnis des ewigen Lebens zu verraten. Ich spürte, wie es versuchte, von meinem Körper Besitz zu ergreifen, wie sein Geist in den meinen schlüpfte... und spürte, dass etwas nicht in Ordnung war, als es statt der erwarteten Todesangst kalte Entschlossenheit vorfand.

Schnell jetzt! Ich ließ mich zu Boden sinken und umfasste den weißen Kristall fest mit beiden Händen. Wer weiß, was geschehen wäre, wenn ich ihn verfehlt hätte. Das Drachenherz nach oben haltend, schrie ich zwei Sätze in der Alten Sprache.

Für einen kurzen Augenblick geschah nichts und ich fürchtete schon, die Kontrolle über meinen Körper verloren zu haben. Dann würde es nicht mehr lange dauern, bis das Böse meinen Geist ausgelöscht hätte. Doch da wurde mir das Böse auch bereits wieder aus den Adern gesaugt. Der Windstoß in der Höhle versiegte und der weiße Kristall in meinen Fingern wurde abrupt eiskalt.

Ich atmete tief durch. Es war mir gelungen! Ich hatte ihn gefangen!

Nun galt es nur noch herausausfinden, wie lange er eingesperrt bleiben musste, ehe er mir mitteilen würde, was ich wissen wollte.





\newpage
\section{Der hellsichtige Geisterjäger}

„Tut mir schrecklich leid, dass du dich meinetwegen so abmühen musst“, entschuldigte sich Leander, „Es war bestimmt unangenehm für dich, so weit in den Norden vorzudringen.“

„Macht Ihr Witze?“, entgegnete Jorn, der Agren, begeistert, „Es war grandios, den sagenhaften Baum der Lieder und die Weiten des Rietgrases mit eigenen Augen zu sehen. Könnt Ihr Euch das vorstellen – Unterstände oberhalb der Erde zusammenzubasteln statt in den Gängen unter der Erde zu wohnen? Ach, was erzähle ich da. Natürlich könnt Ihr Euch das vorstellen, wohnt Ihr doch selbst in so einer Hütte. Habt Ihr denn nie Angst, dass die über Euch zusammenbricht? Das wäre ein ziemlich ironisches Ende, meint Ihr nicht auch? Vorsicht, Stolperstein.“

Leander sandte ein Stoßgebet an Mutter Natur, dass sie ihn vor dem Geschwafel seines Begleiters retten möge. Vorsichtig tastete er mit seinen Füßen den großen Stein vor sich aus und suchte Halt. Dann schwang er sich darüber. Etwas knirschte unter seinen Füßen.

Jorn zog scharf Luft ein: „Auweia, ist das etwa... nun, er oder sie ist jedenfalls schon lange tot, da müsst Ihr euch keine Sorgen machen. Vielleicht ein halber Krieger aus dem Unterirdischen Krieg? Ich glaube, in diesem Helm einen mächtigen Schild der Schildzwerge eingraviert zu erkennen. Habt Ihr Euren Fuß vertreten? Könnt Ihr weiterlaufen?“

Leander grummelte etwas, tastete nach Jorns Hand, ergriff sie und lief rasch weiter.

„Nicht so schnell, sonst zieht Ihr uns noch beide in einen Abgrund“, warnte Jorn, „Verzeiht, wir Agren benutzen die Alte Zwergenstraße so selten, dass sie kaum instand gehalten wird. Wir hätten auch gar nicht die Mittel dazu. Kreaturen benutzen sie viel häufiger. Wer nicht gerne zu Wargorfutter wird, hält sich meistens von hier fern.“

Jorn schwieg für einen kurzen Augenblick. Leander vermutete, dass er sich gerade nervös nach Wargors umsah. Leider konnte Leander das Rauschen der Bäume und die Musik der Toten in den Schluchten des Grauen Gebirges nur kurz genießen, ehe sich Jorns Mund erneut öffnete.

„Nicht, dass wir nicht auch mit einem Wargor klarkommen könnten. Also, ich ja nicht. Aber Ihr besitzt ein wenig Kampfgeschick, oder? Ich habe euch mit dem Messer umgehen sehen. Schöne Verzierungen. So etwas finden wir höchstens noch als verrosteten Überrest aus den uralten Kriegen. Ein wenig mehr nach links, Ihr kommt vom Weg ab.“

Leander schwieg, wie er es schon für den Großteil des Wegs getan hatte. Er strich mit seinem Stab über die schön angeordneten Steine, aus denen die Alte Zwergenstraße aufgebaut war. Hoffentlich waren sie bald am Ziel.\bigskip







„Da wären wir, o Seher. Es war mir eine Ehre, Euch bis hierhin zu begleiten. Auch wenn Ihr nicht der gesprächigste wart, habe ich Eure Gesellschaft außerordentlich genossen“, plapperte Jorn. Leander verzerrte seinen Mund zu einem Grinsen und murmelte ein gequältes „Gleichfalls“, ehe er Jorn losließ und mit seinem Stock nach der nächstgelegenen Wand klopfte. Der Klang war dumpf. Ein mit Kräutern und Moos überwachsener Eingang zu einer kleinen Agrenhöhle.

Vorsichtig tastete sich Leander tiefer in die Höhle hinein. Selbst mit gesenktem Kopf stieß der großgewachsene Narkonier immer wieder gegen die erdige Höhlendecke. Warme Luft wehte ihm entgegen und umwehte seine abgemagerten Wangen. So stolperte er weiter auf die Wärmequelle zu, wo er Holz knacken hörte. Ein Lagerfeuer. Dem Geruch nach wurden gerade getrocknete Kräuter geräuchert. Den Geräuschen (oder besser dem fehlenden Geplapper) entnahm Leander, dass sich höchstens eine andere Person hier befand. Er zog tief Luft ein. Der bittere Geruch von Sternkraut trat in seine Nase und ließ ihn aufhusten. Die Ohren gespitzt, glaubte er, rasselnde Atemzüge zu erkennen.

Er räusperte sich.

Der rasche Atemzug, den er nun vernahm, verriet ihm, dass sein Gegenüber ihn nicht erwartet hatte. Die Agren fing sich allerdings rasch und begann mit ihrer rauen Stimme zu sprechen:

„Verzeiht mir, o Seher, ich hatte nicht gedacht, dass Ihr so schnell hierher fändet. Ich bin Rhona, die Stammesälteste dieses Agren-Stammes. Tretet zu mir und nehmt meine Hand, ich führe Euch zu einer passenden Sitzgelegenheit.“

Die Sitzgelegenheit stellte sich als einen ein klein wenig weicheren und klein wenig weniger kalten Teil des kargen Höhlenbodens heraus, aber Leander war nicht verwöhnt und ließ sich nach dem Gewaltmarsch ins Graue Gebirge nur allzu gerne zu Boden sinken. Er ächzte, als die Muskeln in seinen dünnen Beinen zuerst protestierten und sich dann entspannten.

Rhona murmelte etwas vor sich hin und bot dann an: „Ihr müsst hungrig und durstig sein. Wollt ihr frisches Quellwasser? Ich könnte es mit einigen Edelgelb-Blüten aufbessern. Dieses Gemisch hat sich kürzlich diesseits der Korn-Schlucht äußerst beliebt gemacht.“

Leander winkte dankend ab: „Bemüht Euch nicht, Stammesälteste.“

„Bitte, nennt mich Rhona. Die Formalitäten brauchen Euch nicht zu kümmern. Ihr seid schließlich hier, weil wir Eure Hilfe brauchen.“

Endlich kommen wir zur Sache, schmunzelte Leander innerlich. Jorn hatte alles Mögliche verlauten lassen, außer auch nur den kleinsten Hinweis darauf, warum Rhona ihn hatte rufen lassen. Leander hatte natürlich einige Vermutungen angestellt und verschiedenste Steine und Pülverchen eingepackt, dennoch war er aufrichtig gespannt auf Rhonas Gründe.

„Dann erzählt einmal, Rhona. Ich kann mir vorstellen, dass es für Euer Volk einige Probleme bereitete, nach der Befreiung aus der Gefangenschaft dieses finsteren Nekromanten wieder in Eure Höhlen zurückzukehren und eine neue Normalität aufzubauen. Aber das ist nun auch schon wieder mehr als ein halbes Jahrdutzend her. Ich wüsste nicht, wobei gerade ich Euch behilflich sein könnte. Und warum gerade Euch. Kümmert sich sonst nicht der Stammesälteste Grone um den Kontakt mit Menschen von außerhalb des Grauen Gebirges?“

Rhona schwieg. Wahrscheinlich war es nicht angenehm für sie, zurück an die Zeit der Gefangenschaft unter dem finsteren Hademar zu denken. Dann sprach sie: „Um ebenjenen Nekromanten geht es. Und nun ja, Ihr redet nicht mit Grone, sondern mit mir, weil Ihr nicht hier seid auf Gesuch von jemandem außerhalb des Grauen Gebirges. Ihr seid hier, weil Euch jemand vorgeschlagen hat, der eng mit unserem Stamm in Kontakt steht. Wisst Ihr, der Wolfskrieger hat von Euch geschwärmt.“

Das war interessant. Leander hatte bislang mit Orfen nur wenig Kontakt gehabt, und noch weniger von sich selbst preisgegeben. Einige Male hatte Leander Orfen im Wachsamen Wald angetroffen, einmal sogar in Begleitung dieses riesigen schwarzen Wolfs, des Königswolfs. Bei der Gelegenheit hatte Leander versucht, Orfen zum Plaudern über die Herkunft der magischen Fertigkeiten des Königswolfs zu bringen. Er hätte geradesogut auf Zwergenstahl beißen können. Nur zu den Gerüchten, dass im Grauen Gebirge der Geist eines bösartigen Krahders herumspukte, hatte Orfen etwas verlauten lassen. Offenbar hatten ihn die exakten Fragen Leanders beeindruckt. Sah er etwa eine Art Experte für Geisterkunde in ihm? Auf jeden Fall glaubte Leander nun zu wissen, warum er hierhin gerufen wurde. Er kramte in seinem Mantel nach dem weißen Runenstein, der ihm in diesem Fall behilflich sein könnte, und atmete erleichtert auf, als er ihn ertastete.

„Ihr vermutet, dass der Geist des Nekromanten sich noch immer in dieser Welt aufhält?“, riet Leander ins Blaue hinaus.

„Ihr vermutet richtig“, krächzte Rhona, „Es scheint, als hätte ich den Richtigen gerufen. Auf mein Bauchgefühl ist meistens Verlass.“

„Nun, das muss sich erst noch zeigen.“, sagte Leander vorsichtig. Die Gelegenheit, Kontakt mit dem Geist des Nekromanten aufzunehmen, würde ihn natürlich mit großer Freude erfüllen. Er konnte kaum glauben, über was für Wissensschätze der Nekromant verfügen musste. Aber die Helden von Andor hatten den Nekromanten aus der Winterburg verjagt, und diese waren leider meistens sehr gründlich im Auslöschen von Gegnern aller Art.

„Wie kommt ihr auf diese Vermutung, dass der Geist des Nekromanten noch in dieser Sphäre verweile?“, fragte Leander, innig hoffend, dass die Antwort nicht nur in einem Bauchgefühl Rhonas bestünde.

„Ich vermute nicht“, korrigierte ihn Rhona, „Ich weiß es. Finstere Träume plagen mich schon seit Jahren. Seit unserer Befreiung aus der Winterburg, um genau zu sein. Finstere Träume von einer Gestalt mit einer Kapuze. Erst waren es nur Schemen, doch in letzter Zeit kann ich ihre Form deutlich erkennen. Sie wird stärker.“

„Hademar, der Nekromant?“

„Ich weiß es nicht. Aber ich weiß, dass seit einigen Wochen sämtliche Agren, die sich zu tief in die Überreste der Zwergenbauwerke des Grauen Gebirges vorwagen, spurlos verschwunden sind. Fünf an der Zahl waren es. Fünf verschiedene Bauwerke. Und in jeder der fünf Nächte erschien mir im Traum diese Gestalt in einem roten Umhang. Sie hielt einen Stab umklammert, an dem Agrenschädel hingen, nach jedem Vorfall einen mehr. Ihre finstere Kraft nimmt stetig zu.“

Leander legte seinen Kopf schief.

„Fünf kurz aufeinanderfolgende Unfälle sind natürlich Grund zur Sorge, könnten aber reine Werke des Zufalls sein. Und die Träume... besitzt Ihr überhaupt die Gabe des zweiten Gesichts?“

„Ich weiß, wann ich meiner Intuition vertrauen kann. Spottet nicht, o Seher. Ich bitte Euch bloß, dieser Sache nachzugehen. Falls der Nekromant weiterhin hier umgeht, sind wir nicht sicher. Er dürstet danach, unser Volk zu versklaven und einzukerkern. Nicht einmal zu seinem eigenen Vorteil nutzen will er uns, wie er es mit den Arpachen tat. Es ist reine Rache, nach der er strebt. Rache dafür, dass wir einst seinen Bruder unterstützt hatten, als dieser sich von ihm abwandte und ihn in Krahd zurückließ.“

„Ich werde sehen, was ich tun kann“, murmelte Leander enttäuscht. Fünf verschwundene Agren und einige Träume einer Greisin. Dafür hatte sie ihn die lange Reise ins Graue Gebirge antreten lassen? Nicht, dass er die Macht der Träume nicht kannte, aber das Lesen der Sprache der Gegenwart und der Zukunft war eine höchst komplexe Angelegenheit, deren Studium Jahre des Übens bedurfte, ehe man auch nur halbwegs konsistente Resultate liefern konnte. Bauchgefühle waren hier fehl am Platz.

Die alte Reka wäre da anderer Meinung gewesen, erinnerte Leander sich. Die beiden hatten sich schon tiefgehend über Prophezeiungen und Visionen zerstritten. Im Gegensatz zu ihm verließ sich Reka ausschließlich auf Träume und Bauchgefühle. Nun, Leander wusste ja, wer von ihnen beiden stets korrekte Vorhersagen machte und wer sich in seinen Interpretationen oft hatte fehlleiten lassen.

„O Seher?“, fragte Rhona plötzlich mit einem alarmierten Unterton in der Stimme und unterbrach Leanders Gedanken.

Leander gab einen fragenden Laut von sich.

„Mögt Ihr mir mitteilen, warum ihr ein Herz des Nehal mit Euch tragt?“

Leander erschrak und ließ hastig den weißen Runenstein in seiner Manteltasche los. Dieser weiße Edelstein war einer der größten Edelsteine, den die Schildzwerge je in Cavern gefunden hatten, und einer der beiden Edelsteine, die in die Statue zu Ehren des Drachen Nehal am nördlichen Ende der Korn-Schlucht eingelassen worden waren. Sie waren schon vor Jahrhunderten im Innern der Statue bei Nehals Stein an der Stelle platziert worden, wo sich bei Drachen die Herzen befinden.

Die beiden Kristalle waren kürzlich von einem dreisten Dieb namens Ken Dorr aus der Statue gebrochen worden und hatten ihren Weg in den Besitz eines gierigen Handelszwergs, Garz genannt, gefunden, welcher zumindest einen davon für eine horrende Summe an Leander verhökert hatte. Das andere Drachenherz war wahrscheinlich wieder zurück in Cavern gelangt, wo es nun wie ein Augapfel gehütet wurde. Doch im Gegensatz zu den Schildzwergen wusste Leander um die Magie, die in diesem Edelstein aus den Tiefen der Erde konserviert war. Die nötigen Runen, um diese Magie freizusetzen, hatte er schon längst in seine Oberfläche geritzt. Und so hatte er den Runenstein mit ins Graue Gebirge genommen, weil er geahnt hatte, dass er seine Kräfte hier würde nutzen können.

Hatte die Stammesälteste die Magie in seiner Manteltasche gespürt? Oder war auch das nur ein Bauchgefühl gewesen, eine Vermutung, um zu zeigen, dass sie doch über eine gewisse Gabe des zweiten Gesichts verfügte? Leander drehte sich zur Seite und überlegte sich eine angemessene Antwort, da fuhr Rhona bereits fort.

Ihre Stimme hatte einen drohenden Unterton angenommen und zitterte leicht: „Wenn ich mich nicht stark irre, sollte sich dieser Edelstein in einer Statue zu Ehren des ehrenwertesten Drachen der uralten Zeit befinden, nicht im Mantel eines düsteren Sehers aus dem Drachenland. Es scheint mir respektlos gegenüber Kreatok und Nehal, das Mahnmal an die Macht der Drachen und die Fragilität von Bündnissen aus Gründen der Gier zu schänden.“

Leander suchte in seinem Kopf nach besänftigenden Worten, als ihm plötzlich ein beunruhigender Gedanke kam. Er hatte Gerüchte über das Graue Gebirge gehört, Gerüchte, die Orfens Antworten damals hatten zu einer konkreten These verfestigen können. Aus dieser These folgerte Leander nun eine Vermutung für die konkrete Situation, in der er sich befand. Warum die Agren den weißen Edelstein hatte spüren können. Warum er sie so enervierte. Das war keine Wut, das war... Furcht. Wenn Leanders Vermutung der Wahrheit entsprach, dann nahm der Hilferuf der alten Agren plötzlich ganz andere Dimensionen an.

Glücklicherweise gab es einen direkten Weg, seine Vermutung zu widerlegen, sollte sie falsch sein. Leander ignorierte das aufgebrachte Gezeter der Stammesältesten, zog den weißen Edelstein aus seiner Tasche hervor, hielt ihn dorthin, wo er das Rhonas Gesicht vermutete, und sprach einen Satz in der Alten Sprache.

Wie erhofft, unterbrach die Stammesälteste ihre Tirade sofort, heulte auf und stolperte zurück. Leander vernahm hastige Schritte von weiter draußen und eine ihm unbekannte Agren schrie auf: „Rhona! Was ist geschehen?!“

„Raus mit dir!“, zischte Rhona zurück, „Ich bin bloß hingefallen. Wenn ihr Jungspunde Euch nur ein einziges Mal nicht zu viel Sorgen um mich machen würdet...“

Langsamere, schlurfende Schritte zeugten davon, wie die fremde Agren den Raum beschämt wieder verließ.

Leander konnte es nicht lassen, den Besserwisser zu spielen, beugte sich vor und flüsterte zu Rhona: „Das Drachenherz ist in meinen Händen besser aufbewahrt als im Körper einer Steinstatue. Habt Ihr mich nicht rufen lassen, um den alten Zwergenruinen einen Geist auszutreiben? Dieser Edelstein wurde mit den passenden Runen versehen und ist nun in der Lage, eine nicht mehr in ihrem eigenen Körper verweilende Seele zu greifen und in seinem kristallenen Innern einzusperren. So kann man fremde Geister fangen. Alles, was es dafür braucht, sind die richtigen Worte in der Alten Sprache. Aber das wisst Ihr alles natürlich bereits, sonst wärt Ihr nicht so hektisch zurückgewichen. Euer Glück, dass ich nicht die richtigen Worte gesprochen habe.“

Die Reaktion der Stammesältesten hatte Leander mit ziemlicher Sicherheit erfüllt, dass er gar nicht mit der echten Stammesältesten sprach. Nur mit einem Geist, der ihren Körper führte. Und da er abgesehen von Hademar nur von einem möglichen Geist im Grauen Gebirge vernommen hatte, wagte er einen Tipp abzugeben, mit wem er hier sprach: „Nomion. Der erste Krahder. Es ist mir eine Ehre, Eure Bekanntschaft zu machen.“

Sein Gegenüber blieb still, was Leander großzügig als Bestätigung seiner Vermutung interpretierte. So dachte er laut weiter:

„Was könnte den Geist des Entdeckers der Dunklen Hexerei dazu bewegt haben, den Körper dieser Stammesältesten einzunehmen und nach der vermutlich einzigen Person im Lande zu rufen, die ihn tatsächlich erkennen könnte? Warum sich dem durchaus beträchtlichen Risiko aussetzen, von mir eingefangen zu werden? Hochmut könnte es natürlich sein, Hochmut, wie man es von den Riesen kennt. Ihr hieltet es für unwahrscheinlich, dass ich Eins und Eins zusammenzählen könnte, und wenn ihr nicht so direkt nach dem Drachenherz gefragt hättet, hätte es wohl selbst mich ein bisschen länger gekostet, den richtigen Schluss zu ziehen.“

Die Stammesälteste gab ein Schnauben von sich. Leander fuhr genüsslich fort: „Ihr seid dieses Risiko willentlich eingegangen. Ihr braucht also etwas von mir. Ihr riefet mich hierher und gabt mir den Auftrag, die Ruinen der alten Zwergenbauwerke nach einem Unheil abzusuchen. Da liegt der Schluss nahe, dass in diesen Ruinen tatsächlich etwas umgeht. Etwas, oder jemand. Jemand, vor dem sich selbst der große Nomion fürchtet?“

Leander legte seinen Kopf schief. Da endlich meldete sich die Stammesälteste wieder zu Wort. Ihre Stimme klang plötzlich finster, als würde sie jedes Wort wütend durchkauen und dann erst ausspucken: „Wie wagst du es, so mit mir zu reden? Ich bin Nomion, der Hexer aus Krahd! Der Meister des Urtrolls! Ich gebiete über eine Naturgewalt und ich fürchte mich vor niemandem!“

Leander klatschte in seine Hände und verstaute den weißen Runenstein wieder in seiner Manteltasche. „Warum nicht gleich so?“, fragte er fröhlich. Ohne auf eine Antwort zu warten, führte er seinen Gedankenschluss zu Ende:

„Nomion, der Meister des Urtrolls, musste sich tatsächlich vor niemandem fürchten, als der Urtroll noch das mächtigste Wesen im Grauen Gebirge war. Aber das war einmal. Der Urtroll ist ein Schatten seines früheren Selbst. Jahrhunderte hat er geschlafen, und als er das letzte Mal geweckt wurde, war es nicht von Euch. Er wurde von den Helden von Andor besiegt und von einer Fee kontrolliert, aber auch diese sind es nicht, vor denen Ihr Euch fürchtet. Nein, ihr fürchtet Euch vor demjenigen, der es vermochte, Euren Urtroll mit einem einzigen grellen Strahl aus purer Magie zu vertreiben.“

„Nomion fürchtet sich vor niemandem!“, keifte Nomion erneut auf. Leander musste ein Glucksen unterdrücken. War dies etwa wirklich der bösartige Herrscher, der einen Kult von Krahder-Hexern gegründet und die Drachen beinahe vollständig ausgelöscht hatte? Dessen Körper nur von Tarok selbst hatte vernichtet werden können und dessen Geist sich selbst danach noch an diese Welt hatte klammern können? Der Zahn der Zeit schien auch an ihm genagt zu haben.

„Sagt mir, Nomion“, fragte Leander nun, „Ist es Hademar von Krahd gelungen, die Lehren der Dunklen Hexer mit geheimen Wissen aus der Akademie von Hadria zu kombinieren? Hat der Nekromant, gestärkt durch die Drachenmagie, es geschafft, dem Tod ein noch größeres Schnippchen zu schlagen, als Ihr es vermochtet? Vermutet Ihr etwa, dass sein Geist dem Euren überlegen ist?“

„Hademars Körper wurde von den Helden von Andor vernichtet, so wie Tarok den meinen zu Asche vergehen ließ“, knurrte Nomion, „Nicht aber sein Geist. Hademars Geist streift immer noch in diesen Landen umher, setzt sich langsam zusammen und wird mit jedem vergangenen Tag stärker. Meine Kräfte sind nicht mehr, was sie einst waren, und so ließ ich dich rufen, weil ich dir die Gelegenheit geben will, Hademar ein für alle Mal auszulöschen.“

Ein Feigling ist Nomion, wie alle Krahder, dachte Leander im Geheimen. Äußerlich sprach er: „Eine spannende Gelegenheit ist das zweifelsohne, o Erster der Krahder. Doch verratet mir dies: Warum sollte ich sie ergreifen?“

Nomion quetschte seinen nächsten Satz hervor und Leander hörte förmlich, wie er dabei die Lippen der Stammesältesten zu einem hämischen Grinsen verzerrte: „Mach mir nichts vor, Leander, Seher von Narkon. Ich habe von dir gehört. Du verzehrst dich doch förmlich nach vergangenem Wissen, und wer auf dieser Welt außer Hademar kennt sich schon sowohl in der unsrigen als auch in der nordischen Magie aus? Sein Fang wäre eine Goldgrube für dich. Hopp, hopp, Silberländler. Dein Ziel hast du doch bereits festgelegt, sobald ich dir verraten habe, dass Hademars Geist noch in der Nähe verweilt und eingefangen werden kann. Möge die Jagd beginnen.“

Leander trommelte mit seinen Fingern auf seinem Holzstab herum. Nomion hatte recht, Leander würde sich nicht im Traum die Gelegenheit entgehen lassen, Hademars Geist zu fangen und seine Kenntnisse auszupressen. Doch mochte er es nicht im Geringsten, ausgespielt zu werden. Normalerweise war Leander es, der andere Personen zu seinen eigenen Zwecken manipulierte. Dass Nomion den Spieß nun umdrehte, gefiel ihm nicht im Geringsten. Leanders eigenen Zwecke bestanden oft darin, heldenhafte Gestalten an seiner Stelle in den gefährlichen Norden Silberlands zu schicken. Heldenhafte Gestalten wie den melancholischen Schwarzen Barden Orril. Oder wie den Feuerkrieger Trieest aus dem fernen Danwar. Wenn damals diese Wassermagierin Jarid nur nicht dazwischen gekommen wäre...

Leander musste an Rhona denken. Wie lange hatte sie gemeinsam mit dem restlichen Volk der Agren unter dem Joch Hademars leiden müssen, nur um nun unter Nomions Kontrolle zu stehen? War sie wohl wach? Konnte sie fühlen, was um sie vorging? Beim Gedanken an ihr Leiden wallten Schmerzen in Leander auf.

Schnell schüttelte Leander diesen Gedankengang beiseite. Die Stammesälteste bedeutete ihm nichts. Er hatte selbst schon so viele Personen kontrolliert, wenn auch nie so direkt wie Nomion. Die Frage, die ihn vielmehr beschäftigen sollte, war, wann und wie öffentlich Nomion ihren Körper verlassen würde. Wenn der Agrenstamm Wind davon bekäme, dass Leander gewusst hatte, was mit ihrer Stammesältesten vorgegangen war, und nichts dagegen getan hätte, dann würden die Agren wohl kaum mehr so freundlich mit ihm umgehen. War das ein verkraftbarer Verlust?

Während Leanders Gedanken rasten und kalkulierten, wandte er sich wieder an Nomion und erwiderte: „Hademars Kenntnisse über die Dunkle Hexerei umfassen doch höchstens das wenige, was er als Sklave in Eurem Tempel aufgeschnappt hatte. Sie können unmöglich so tief gehen wie die Euren, der diesen Tempel doch überhaupt erst errichten ließ. Gibt es einen Grund, warum ich Euch nicht an Ort und Stelle im Drachenherz gefangen nehmen und Euer Wissen auspressen sollte?“

Urplötzlich wurde es kalt um Leander. Jedes Haar auf seinem Körper stellte sich auf. Er hörte die Stammesälteste röchelnd husten und spürte die Gegenwart von etwas Fremden. Eine Macht, die groß genug war, um einen jeden Geist zu versklaven. Leander nahm wahr, wie \textit{Etwas} in seinen Geist eindrang und seine Gedanken in Eiseskälte tauchte. Er stolperte alarmiert zurück und stieß mit seinen Kopf gegen die Höhlendecke.

Sein Körper fing sich von selbst. Von Horror erfüllt, spürte Leander seine Zunge in seiner Mundhöhle umherwandern und seine Zähne entlanggleiten. Seine eigenen Lippen öffneten sich und er hörte seinen eigenen Mund krächzend sprechen: „Drohst du mir etwa, Leander von Narkon?!“

Leander versuchte zu antworten, versuchte, zurückzuweichen, versuchte, in seiner Tasche nach dem Drachenherz zu langen. Alles erfolglos, sein Köper gehorchte ihm nicht mehr.

Wie ein Fremder in seinem eigenen Kopf vernahm er, wie sein Mund sich erneut verdrehte und sprach: „Na los, versuch schon, mich im Edelstein einzufangen.“

Ein grässliches Lachen entsprang Leanders Kehle, als Nomion fortfuhr, nun aus Leanders Mund statt aus Rhonas: „Es kann nur eine Seele auf einmal im Drachenherz gefangen sein. Hademar wird stetig stärker und muss möglichst bald gebändigt werden. Seine Kenntnisse werden dich für Jahre beschäftigen. Das muss deine Priorität sein. Später, wenn du dich an ihm ergötzt hast, magst du ihn vernichten und dich wieder auf die Suche nach mir machen. Aber wisse, dass du mich nur dann finden wirst, wenn ich es wünsche. Denn alles geschieht, wie Nomion gebietet!“

Leanders Magen drehte sich auf den Kopf, als Nomions Kälte seinen Körper verließ. Röchelnd stürzte Leander zu Boden und hustete sich die fremde Seele aus dem Leib. Ein letztes Lachen vernahm er noch („Ich hoffe für dich, dass du dich bei Hademar geschickter anstellst als gerade eben!“), dann huschte ein Windhauch an ihm vorbei und Nomions Präsenz war nicht mehr da.

Leander und die Stammesälteste waren beide alleine und lagen schwer atmend am Höhlenboden. Leander rührte sich als erster und fragte: „Stammesälteste! Rhona! Seid Ihr verletzt? Geht es Euch gut?“

Die Stammesälteste antwortete nicht in Worten, sondern stieß einen langanhaltenden Schrei aus, der all ihre Qualen ausdrücke, die sie nicht auszuformulieren vermochte.

Leander wich schockiert zurück. Wie lange war sie die Gefangene Nomions gewesen?

Erneut hörte Leander hektische Schritte die Höhle betreten und eine fremde Agrenstimme fragte: „Rhona, Rhona, so sprecht mit mir! Was ist los?“

Die Stammesälteste schluchzte auf und schrie: „Was... was ist geschehen? Welcher Dämon... bei der Mutter der Erde, wer war das?!“

Leander wandte sich an die Neuankömmlinge und sprach: „Rhona wurde von einem bösartigen Geist erfüllt, doch dieser ist weitergezogen. Er kann ihr nicht weiter schaden.“

Weitere Schritte, dann wurde Leander plötzlich am Kragen gepackt: „Was hast du mit Rhona gemacht, du Finsterling?!“

„Wie ich bereits sagte, sie wurde einem finsteren Geist kontrolliert. Ich hatte nichts damit zu tun, ich habe sie vielmehr befreit.“

Agren waren genau wie die Mehrzahl der Menschen, Zwerge und Taren von Emotionen erfüllt, die ihnen das Zusammenleben vereinfachten und sie offener für Leanders Manipulationen machten, aber leider auch oft von rationalen Handlungen abhielt. Das musste Leander erneut erleben, als ihn plötzlich etwas Spitzes in die Seite stach. Verächtlich schlug er der Agren die Klinge aus der Hand, da schlug etwas anderes, hölzernes, feste gegen sein Schienbein. Leander fluchte auf. Den Schmerz zu ignorieren vermochte er natürlich, aber sein Bein war kein Baumstamm und gab unter dem Schlag nach. Leander ging zu Boden. Sein Kopf schlug gegen den Höhlenboden, der nun doch nicht mehr so weich wirkte, wie als Leander sich zuerst darauf niedergelassen hatte.

Alles wurde farbig. Leander sah eine Vision der Winterburg sich um ihn herum aus dem Boden erheben und konnte durch die Mauersteine hindurch den Körper der Stammesältesten erkennen, der in einer kleinen Zelle saß. Dann flog die Gittertür auf und eine seltsame Schattengestalt trat ein, die sich stetig deformierte und zwischen der Silhouette eines Menschen und der eines Bären hin- und herzuwechseln schien. Der Mund der Stammesältesten verzog sich zu einem Lächeln und alles wurde wieder schwarz.

„Der Mensch im Bären und der Bär im Menschen wird kommen, um Euch zu retten“, flüsterte Leander, ohne diese Worte gewählt zu haben.

„Zurück, tretet zurück!“, erklang die herrische Stimme Rhonas, „Er spricht die Wahrheit, nicht er hat mir etwas angetan.“

Leander bemerkte erleichtert, dass ihn keine Stiche und Schläge mehr trafen. Vollkommene Stille erfüllte die Höhle, oder zumindest derart prominente Stille, wie sie in der Präsenz mehrerer Personen auftreten konnte. Manchmal war Leander froh darüber, dass seine Augen aufglühten, sobald er eine Vision empfing. Das erlaubte seinem Umfeld, zu erkennen, dass er nicht einfach Lügenmärchen erfand.

Als die Stimme der Stammesältesten erneut erklang, war sie ganz nahe an Leander getreten und hauchte ihm förmlich ins Gesicht: „Was... was habt Ihr soeben gesagt? Was habt ihr gesehen, Seher?“

„Es tut mir so leid“, antwortete Leander, „Ihr werdet noch ein drittes Mal in Gefangenschaft geraten. Ich habe es gesehen. Doch fürchtet Euch nicht. Der Mensch im Bären und der Bär im Menschen wird kommen, um Euch zu retten. Und danach werdet Ihr den Rest Eures Lebens in Freiheit genießen können.“

Leander richtete sich abrupt auf und richtete seinen Mantel. Er tastete sein Schienbein ab. Nichts schien gebrochen und den Stich in seine Seite fühlte er nicht einmal mehr.

„Stab!“, verlangte er herrisch. Von seiner Rechten wurde ihm sein Stab in die Hand gedrückt.

„Aus dem Weg!“, brachte er noch heraus, dann stolperte Leander mehr aus der Höhle, als dass er schritt.\bigskip







Erst als Leander sich einige Minuten von der Höhle der Stammesältesten entfernt hatte, hielt er inne und verschnaufte. Er rückte seine Augenbinde gerade und fühlte überrascht Tränen seine Wangen runterrollen. Die Vision hatte ihn stärker getroffen, als er angenommen hatte. Leander sank zu Boden.

Die Wahrheit war, dass dies das erste Mal seit Langem gewesen war, dass ihm endlich wieder eine Vision erschienen war. Natürlich hatte er gewusst, dass er seine Gabe der Voraussicht nicht einfach so verloren hatte. Dennoch war in den letzten Monaten ein immer heftigerer Samen der Furcht und des Zweifels in ihm gewachsen. Die Bestätigung, dass er noch immer das Zeug zum Seher hatte, erfüllte ihn mit Erleichterung und Euphorie. Die Hoffnung, seinen Bruder doch noch retten zu können, erwachte wieder in ihm.

Leander erinnerte sich daran, als wäre es gestern gewesen, was er sich damals davon erhofft hatte, die Sprache der Zukunft zu erlernen: Ein Gefühl für die Zukunft, eine Gewissheit, was ihn erwarten würde. Seinen Tod würde er voraussehen, und zwar so weit herausgezögert, wie nur irgendwie möglich, denn wenn es ein vermeidbarer Tod wäre, so könnte er ihn ja vermeiden und hätte ihn gar nicht erst gesehen. Allerlei geniale Ideen und Pläne aus der Zukunft würde er empfangen, und funktionieren würden sie, weil, wenn sie noch verbessert werden könnten, so hätte er sie ja bereits verbessert vorhergesehen. Ja, damals hatte Leander noch gedacht, dass seine seherischen Fähigkeiten ihn zu einer Art Gottheit machen würden, fast allwissend, fast unaufhaltsam.

Doch es war nicht dazu gekommen. Alles, was Leander vernehmen konnte, waren kurze Einblicke in die Zukunft, Schemen, Gestalten, Wortfetzen. Auch wenn er die Sprache der Zukunft erlernt hatte, war ihre Stimme oft leise und schwer verständlich. Und am Schlimmsten war, dass er nicht einmal kontrollieren konnte, wann ihn eine Vision ergriff. Zufällig waren sie nicht, dafür tauchten sie zu oft in passenden Situationen auf. Ob eine bewusste höhere Macht dahintersteckte, die ihm diese Visionen zuschickte? Diese würde ihn dadurch gut manipulieren können. Das war ein Bild, das Leander nicht gefiel. Und wer würde ihm schon nur so spärliches Wissen zukommen lassen wollen und vor allem zu welchem Zweck? Andererseits hatte Leander auch keine andere eindeutige Gesetzmäßigkeit hinter den Visionen erkennen können, allerhöchstens Trends. In ereignisreichen Jahren tendierten mehr Bilder der Zukunft dazu, ihn zu erreichen. Insbesondere wenn Leander selbst betroffen oder in Gefahr war. Dem gegenübergestellt hatte Leander in den letzten Monaten gleich zwei brenzlige Konfrontationen mit marodierenden Gorlots erlebt, ohne auch nur eine Vision zu empfangen. Die Sachlage blieb unergründlich.

Auf jeden Fall empfand Leander ungemeine, ja, gar überwältigende Freude darüber, nun endlich wieder einen Blick in der Zukunft empfangen zu haben, auch wenn dieser sich nur um die unwichtige Stammesälteste gedreht hatte.

Leander richtete sich auf und versuchte, diese Gefühle und diese Gedanken an Vergangenes zu unterdrücken.

Es gab im Moment Wichtigeres zu tun.

Er hatte einen Nekromanten zu fangen.











\newpage
\section{Nervige Hexen, zwei an der Zahl}


Leander tastete den Boden nach einer Schlafgelegenheit ab. Sein Stock stieß auf Knochen, Knochen und noch mehr Knochen. Schädel, Rippen und Beine von Menschen, Zwergen und Schafen zugleich übersäten die Knochengrube. Hier und dort konnte man sogar den magischen Fußknochen eines waschechten Drachen finden.

Eigentlich sollte Leander in Hochstimmung sein. Er hatte soeben (im Prolog) geschafft, Hademars stetig stärker werdenden Geist zu einem Blutsteinaltar der Zwerge unter dem Grauen Gebirge zu rufen und im Licht des roten Mondes ins Drachenherz zu sperren. Nun leuchtete dieser einst strahlend weiße Edelstein in einem tiefdunklen Rot. Zumindest hatte ihm Jorn der Agren wortreich von der satten Farbe geschwärmt. Jorn hatte zudem unter der steinernen Oberfläche hin und wieder dunkle Schwaden zu erkennen geglaubt, welche willkürlich im Innern des Drachenherzes umherwaberten. Hademars Geist musste sich wohl erst an seinen neuen kristallenen Körper akklimatisieren und herausfinden, wie er seine schwarze Seele im Kristall verformen konnte, um mit der Außenwelt zu kommunizieren. Nun, Leander würde ihn darin unterstützen, sobald er ihn erst einmal zu seiner Hütte im Wachsamen Wald zurückgebracht hatte.

Siegestrunken war Leander zurück nach Andor aufgebrochen, diesmal ohne die Begleitung Jorns. Doch hier, in der Knochengrube, wo er sein Nachtlager aufschlagen würde, konnte er einfach nicht fröhlich bleiben. Leander trauerte. Er trauerte um die Drachen, deren kollektives Wissen, deren Magie, deren Geschichten und deren Gedanken ein für alle Mal ausgelöscht worden waren, als diese elenden Helden von Andor den letzten ihrer Art erschlagen hatten.

Leander erinnerte sich nur allzu gut daran, wie Schreie durch den Wachsamen Wald gedrungen waren und wie ein Geruch nach Rauch und Schwefel das ganze Land erfüllt hatte. Das war einer der Momente gewesen, wo sich Leander seine Sehkraft zurück gewünscht hatte. Taroks Flug über das Land hätte er nur zu gerne wahrgenommen. Wie die riesige Echse ihren Schatten über das Rietland warf, wie sie vom Himmel stieß... aber dann waren ihre Flügel von Pfeilen durchlöchert, ihre Schnauze von Äxten zerkratzt und ihre Herzen von Schwertern durchstoßen wurden. Barbaren waren sie doch, diese Helden! Und nicht nur Fenn.

Leander hätte einen friedlichen Weg gefunden, Tarok zu besänftigen. Den Zugang zu Krahal, all dieses Wissen, verloren für immer. Was für eine Schande!

So trauerte Leander lange um Tarok, ehe er endlich einschlief.\bigskip







Leander erwachte aufgrund der Kälte. Kälte war keine Besonderheit im Grauen Gebirge und für Leander wäre es nicht das erste Mal gewesen, dass sein pelziger Umgang des nachts zur Seite gerutscht wäre und ihm die Wärme seines Körpers verwehrt hätte.

Doch so kurz nach seinen Begegnungen mit den Geistern Nomions und Hademars war sich Leander umso stärker bewusst, wie sich eine unnatürliche Kälte anfühlte. Und die Kälte, die gerade Leanders Gliedmaßen schlottern ließ, war definitiv keinen natürlichen Ursprungs.

Leanders Vorgehen war rasch und methodisch.

Sein erster Gedanke galt Hademar, falls dieser es geschafft hatte, dem Drachenherz zu entrinnen. So tastete Leander in seiner Manteltasche nach dem Runenstein und bestätigte erleichtert dessen Vollständigkeit und Unversehrtheit.

Leanders zweiter Gedanke galt Nomion. War der Geist des ersten Krahders zurückgekehrt, um dem eingekerkerten Hademar ein für alle Mal den Garaus zu machen? Ein perfider Plan, der einem hinterlistigen Krahder alle Ehre machen würde, doch würde Nomion es wohl kaum riskieren, dass Leander seine Ankunft bemerkte. Leander konnte Hademar jederzeit aus dem Edelstein befreien und dann würde Nomion einem wahnsinnigen Nekromanten gegenüberstehen, der die Macht der Drachenmagie für sich nutzen konnte.

Nein, diese unnatürliche Kälte schien weder von Hademar noch Nomion zu stammen. Leander ging im Geiste sämtliche Begegnungen durch, die ihn an die jetzige erinnerten, und fand nach einem kurzen Augenblick einen Treffer. Genau diese Form von schleichender Kälte, die ihn an eine mondlose Nacht und an ein schattiges Tal erinnerte, diese Form von Kälte hatte er schon einmal gespürt. Aber das war lange her gewesen. Und Leander hätte schwören können, dass ihre Urheberin bereits vor über einem halben Jahrdutzend von diesen elenden Helden von Andor vernichtet worden war, tatsächlich erst wenige Tage vor Taroks Tod.

Schlaftrunken richtete sich Leander auf, griff nach seinem Holzstab und rief: „Alte Freundin, was schleichst du dich so an? So zeige dich doch! Wie hast du es denn geschafft, dem Getümmel bei der Rietburg zu entkommen?“

Mit Schrecken erkannte Leander, dass etwas Eiskaltes seine Füße umspülte und langsam seine Beine hochstieg.

„Lass das!“, fauchte Leander mit einer ein bisschen zu hohen Stimme, um einschüchternd zu wirken, „Ich weiß, du bist die Dunkelheit in Person und willst alles verschlingen, was sich dir entgegen stellt, aber ich kann dir so viel nützlicher sein, wenn du mich fürs Erste verschonst. Wir hatten eine solche Konversation schon einmal! Ich kam dich in deinem Verlies an der Rietburg besuchen. Hast du denn alles vergessen in der Zwischenzeit?“

Leander hatte sie tatsächlich im Verlies der Rietburg kennengelernt, als er sie in der Rolle eines mürrischen Gelehrten aus dem Wachsamen Wald aufgesucht hatte, um Kenntnisse über ihre Kräfte zu erlangen. Sie war seit den Trollkriegen dort angekettet und halb wahnsinnig gewesen, doch hatte sie sich ihm gegenüber aus ihm unverständlichen Gründen vollkommen freundlich verhalten. Nun wirkte sie überhaupt nicht mehr freundlich, als eine flüsternde und doch durchdringende Stimme erklang: „Wir kennen uns nicht, Fremder, doch sei getröstet: Sobald ich mir deine Seele einverleibt haben werden, werde ich dich kennen.“

Ja, dann ist es zu spät, dachte Leander trocken.

„Shan!“, rief er aus, „Ich will dir nicht wehtun, aber ich werde es tun, wenn ich muss.“

Shan, die Schattenhexe, lachte nur leise vor sich hin, während ihre Schattententakel sich noch stärker um Leanders Beine wanden.

Genug war genug. Leander konnte sich zwar nicht erklären, warum Shan sich nicht an ihn erinnerte, aber ihm war klar, warum sie hier war. Die Schattenhexe dürstete es nach Seelen und Hademars von Magie nur so strotzende Seele war im Drachenherz auf engsten Raum gequetscht. Eine so dichte, so mächtige Seele musste auf die hungrige Shan wie ein Leuchtfeuer wirken, das sie magisch anzog. Praktischerweise war der Kristall, dessen Inhalt Shan derart anzog, auch einer der einzigen effektiven Schutze gegen Shans Präsenz: Ein Edelstein Caverns, in welchem das Licht der Welt gespeichert war.

Natürlich förderte Hademars Präsenz mit all ihrer Dunklen Hexerei und Dunklen Magie das Licht des Drachenherzens nicht im Geringsten. Dennoch hoffte Leander, dass genug von der ehemaligen Reinheit des weißen Edelsteins übrig war, um Shan zu vertreiben.

So klaubte Leander zitternd das erkaltete Drachenherz aus seiner Manteltasche und hielt es Shans Stimme entgegen. Prompt erklang ein durchdringendes Heulen, als Shan zischte:

„Es brennt! Steck es weg! Weg!“

„Weiche zurück, und ich stecke den Stein weg, Shan!“, rief Leander aus, „Ich weiß, dass diese Seele unglaublich kräftig ist, aber du kannst sie nicht haben. Zumindest noch nicht.“

Shan heulte erneut auf, aber ihre Tentakel ließen endlich Leanders Beine los und der Seher stolperte zurück. Fast war er froh, dass er nicht sehen konnte, wie nahe Shans Schatten von allen Richtungen an ihn herandrangen, als er das Drachenherz in seinem Mantel verbarg, aber immer noch in seinem Griff behielt.

„Verflucht seist du, Fremder, verflucht! Mögest du nie wieder Frieden finden, bis...“, wisperte Shan.

Leander lachte auf. Er selbst hatte oft genug versucht, Seekönig Varatan zu verfluchen, als dass er sich von Shans Drohungen beeindrucken ließe: „Flüche bedeuten nichts, sofern sie nicht von einer Person mit königlichem Willen gesprochen werden.“

Und Personen mit königlichem Willen hatten oftmals keine Ahnung, was sie mit einem Fluch auslösten. Varatan, der Seekönig, hatte vor langer Zeit einen Fluch über eine ganze Insel verhängt, der seine Bewohner langsam in den Wahnsinn trieb. Unter diesen Bewohnern befand sich auch Callem, Leanders Bruder. Und Leander hatte seither immer wieder heldenhafte Gestalten nach Silberland geschickt, ohne dass auch nur eine davon hätte herausfinden können, wie man den Fluch brechen konnte. Stattdessen waren sie alle ihm verfallen.

„Ist das so?“, fragte Shan bedeutungsvoll, „Ist es so, dass nur ein willensstarker König einen Fluch verhängen kann? Wie kommt es dann, dass auf mir ein Fluch der Dunkelheit liegt? Die Dunkelheit selbst hat doch keinen eigenen Willen. Oder etwa doch?“

„Manchmal ist es besser zu schweigen, als etwas zu behaupten, von dem man nichts versteht“, stieß Leander müde hervor, „Ziehe von hinnen, Shan! Suche dir eine andere Beute und lass uns in Ruhe. Oder soll ich den Runenstein erneut hervorholen?“

„Es schmerzt, es schmerzt so sehr, diese Seele ziehen zu lassen, und es schmerzt, in ihrer Nähe zu bleiben“, wimmerte Shan, ehe sie fortfuhr: „Hüte deine Zunge, Fremder! Ich verstehe mehr von Flüchen, als du meinen magst! Ich kenne gleich mehrere Verfluchte, und an einem davon zieht ein Fluch, deren Signatur dem ganz ähnlich ist, der schwach an dir haftet. Vielleicht gar identisch.“

Leander zuckte zurück. Er wusste, dass er Varatans Fluch erschreckend nahe gekommen war, damals, als er nach Varatans Angriff auf die Schwarze Kogge die Silberberge erklommen hatte, um Callem zu suchen. Wie aus dem Nichts war das durchsichtige, verfallene Gesicht Varatans vor ihm aufgetaucht und hatte ihn angestarrt. Die Manifestation des Fluchs war jedoch nicht näher gekommen. Inzwischen wusste Leander, dass die Silberberge das Ende des Reviers des Fluchs markierten und dieser sie nicht überqueren konnte. Doch hatte Leander gedacht, dass diese verstörende Begegnung nicht gereicht hatte, um von ihm markiert zu werden.

Er schüttelte seinen Kopf. Das war egal. Das war alles egal. Wichtig war doch nur...

„Beweis es“, flüsterte Leander verschwörerisch, „Zeig mir, dass du dich auskennst, o Shan. Sage mir, wie man Varatans Fluch brechen kann.“

Shan lachte auf: „Das ist doch einfach. Töte Varatan! Dann löst sich der Fluch.“

Leander lachte auf: „Wenn es doch nur so einfach wäre! Varatan ist bereits vor langer Zeit umgekommen.“

In Gedanken ergänzte Leander, dass er sonst ja schon längstens einen Weg gefunden hätte, Varatan persönlich davon zu überzeugen, den Fluch aufzuheben.

Dann fuhr er fort: „Der Fluch wird auch nach seinem Tod weiterhin aufrecht erhalten. Kann das sein?“

„Ja, natürlich“, zischte Shan in aller Seelenruhe, als wäre es das Selbstverständlichste auf der Welt, „Die Nachkommen deines Fluch-Urhebers nähren diesen ebenfalls – sofern sie die nötige Willenskraft dazu besitzen. Varatans Willenskraft lebt in seiner Blutlinie weiter.“

Leander stutzte. Konnte die Lösung tatsächlich so einfach sein? Musste er nichts tun, als Varatans Nachkommen zu vernichten? Er fluchte innerlich auf beim Gedanken, dass er so nahe daran gewesen wäre.

Seine Gedanken flogen zurück in die Zeit seiner Jugend.

Sein Bruder Callem war früh von zu Hause aufgebrochen. Voller Tatendrang und einer unbeschreiblichen Gier nach Ruhm. Leander hingegen war zu Hause geblieben. Hatte lesen und ein wenig die Kunst des Heilens erlernt. Dann hatte er Interesse an Sprachen entwickelt. So hatte er die umliegenden Inseln besucht und ihre Sprachen erlernt. Überrascht hatte er damals festgestellt, dass von der Existenz seiner Heimatinsel auf den restlichen Nebelinseln niemand wusste.

In Varatanien hatte Leander sich niedergelassen und zu einem Physikus ausbilden lassen. So hatte er seine Kenntnisse mehren können und gleichzeitig den Kontakt mit Callem aufrecht erhalten, wenn dieser nicht gerade im Auftrag des Seekönigs das Hadrische Meers befuhr.

Doch dann war alles schief gelaufen. Callems Geist war von Kenvilar, der tückischsten Macht des Hadrischen Meeres, vergiftet worden. Er hatte sein Schiff vernichtet und eigenhändig den Großteil seiner Mannschaft getötet. Unter Kenvilars Führung hatte er eine neue Besatzung versammelt, die Besatzung der gefürchteten Schwarzen Kogge. Und gemeinsam waren diese Unruhestifter durch das Hadrische Meer gezogen, hatten Handelsschiffe überfallen und Waren geraubt.

Aufgrund seiner brüderlichen Verbindung mit Callem war Leander auf den Nebelinseln in Verruf geraten, sobald herausgekommen war, dass Callem noch lebte, in den Bann des Bösen gelangt war und nun nach der Krone der Nordmeere strebte. Daraufhin war Leander auf einem Schiff gen Hadria gereist. Er hatte zwar nie ein Talent für Zaubersprüche und mächtigeren Hokuspokus gezeigt, doch das Wissen der Akademie von Hadria und die Gespräche mit dem Gelehrten Hombudt, dem Zauberer des Raums im Turm der Magie, waren für ihn Grund genug gewesen, sich eine Zeit lang in Nordgard niederzulassen.

So hatte Leander nur durch Hörensagen erfahren, wie Varatan die Schwarze Kogge bis in ihr Versteck verfolgt hatte, Leanders und Callems Heimatinsel. Wie Varatan die Insel Narkon getauft und seinen unsäglichen Fluch darüber verhängt hatte, ungeachtet der restlichen Bewohner der Insel.

Hingegen hatte Leander direkt miterlebt, wie Varatan Jahre später in den Norden gereist war. Er hatte den grünlichen Rauch aus der Hadrischen Unterwelt aufsteigen sehen, als Orweyn mit einem gewöhnlichen Falken und drei gewöhnlichen Waffen in den Horun hinabgestiegen und mit dem riesigen Falken von Yra und drei magischen Waffen zurückgekehrt war. Er hatte aus sicherer Entfernung durch ein Fernrohr beobachtet, wie Varatan mit den drei magischen Waffen stolz die Mächte des Meeres herausgefordert hatte und seine Flotte von den Mächten vernichtend geschlagen worden war. Ja, Leander hatte miterlebt, wie Varatan nach Orweyns Opfer – oder besser gesagt Orweyns Massaker an der Zauberergemeinschaft – seine Königswürde abgelegt hatte. Leider hatte dies nicht gereicht, um seinen Fluch über Narkon aufzuheben.

Leander war auch dabei gewesen, als Varkmar, Sohn des Varatan, unzufriedene Zauberer in Nordgard um sich geschart und den Zaubererorden des Feuers gegründet hatte. Varkmar selbst hatte ihm den entscheidenden Tipp gegeben, wie Leander die Sprache der Zukunft erlernen könnte. Es war eine Form der Dunklen Magie, die ihn befähigte das zu sehen was nicht ist, sondern erst noch kommt. Doch wie immer hatte die Dunkle Magie ihren Tribut gefordert. In gleichem Maße wie Leanders Fähigkeiten wahrzusagen zugenommen hatte, waren seine Augen für das hier und jetzt schwächer geworden, bis ihr Licht ganz erloschen war.

So nahe war Leander an Varkmar gewesen! Eine einzige Phiole schwarzen Vypera-Gifts, ein kleiner Stich mit einer Nadel, und Varatans Fluch wäre gebrochen gewesen.

Doch stattdessen hatte sich Leander, enttäuscht über die vagen Visionen des erst noch Bevorstehenden und bestürzt über den Verlust seines Augenlichts, aus der kalten und klammen Eiswelt Hadrias ins warme Andor zurückgezogen und sich im Wachsamen Wald eine Hütte ergaunert. Die Bewahrer vom Wachsamen Wald betrachteten ihn nun wohl ähnlich wie die Hexe Reka als eigenbrötlerischen mürrischen Weisen, der hin und wieder aufkreuzte, um Wissen zu erfahren oder zum Rechten zu sorgen, und auch nach Jahren kaum einen Tag gealtert wirkte.

So war Hadria aus dem Blickfeld des Sehers geraten, und nun hatte Leander keine Ahnung, wie viele Kinder Varkmar Zeit seines Lebens in die Welt gesetzt hatte. Oder wie viele Kinder diese wiederum gezeugt hatten. Es würde potentiell sehr schwer werden, Varatans Blutlinie zu verfolgen. Er verwarf den Gedanken und wandte sich seiner aktuellen Konfrontation mit der Schattenhexe zu.

„Ich danke dir für diesen Hinweis, Shan“, sagte Leander, „Aber nun muss ich dich endgültig bitten, mich zu verlassen. Meine Seele wirst du nicht kriegen. Und auch das in diesem Edelstein eingesperrte Leben bleibt dir verwehrt. Zumindest heute.“

Ein demonstrativer Griff in seine Manteltasche, wo noch immer das leuchtende Drachenherz steckte, genügte. Die Schattenhexe kreischte ein letztes Mal auf („Die Dunkelheit wird dich bald erneut heimsuchen. Ich werde wiederkehren, wenn du erst einmal ungeschützt bist!“), dann wurde Leander merklich wärmer, als ihre Präsenz diesen Ort verließ. Er rubbelte an seinen Armen und Beinen und versuchte so mehr oder minder erfolgreich, seine Gliedmaßen vom Zittern abzuhalten. Schlussendlich legte er sich sorgsam wieder neben den Knochen der zahlreichen Drachenopfer in der Knochengrube nieder und versuchte, seinen Atem zu beruhigen, während seine Gedanken kreisten.\bigskip







„Das darf doch nicht wahr sein“, fluchte Leander. Nach der langen Reise war er endlich wieder im Wachsamen Wald eingetroffen, hatte sich mit einigen den grünen Radius beschützenden Bewahrern ausgetauscht – keine besonderen Vorkommnisse im Land in seiner Abwesenheit – und sich danach noch auf seine Hütte gefreut, die grau und unscheinbar am Waldrand stand, direkt an der Küste des Hadrischen Meeres. Das Rascheln der Blätter hatte die Luft erfüllte und das entfernte Rauschen der Wellen überdeckt. So hatte Leander gelauscht. Der Wind spielte schon überall sein Lied auf einem anderen Instrument. Die Gerüche des Waldes waren ihm in die Nase gestiegen. Für einen Augenblick hatte ihn das Gefühl erfüllt, zu Hause zu sein.

Doch dann hatte er seine Stirn gerunzelt und gerade eben geflucht. Ein feiner Geruch hatte seine Nase erreicht, der ganz und gar nicht ins restliche Sinnesbild passte. Ein Dampf war es, der Anteile von verdorbener Krallenflechte enthielt, aber auch nach süßen Bärenbeeren roch. Je näher Leander durch den lauschigen Waldweg auf seine Hütte zustapfte, desto lauter wurde das Rauschen des dahinter liegenden Hadrischen Meeres, aber desto stärker wurde auch dieser Geruch nach Kräutern und Rauch.

Kein Zweifel: Jemand kochte gerade in seiner Hütte ein Süppchen! Leander kannte nur eine, die so dreist wäre, ungefragt in sein Heim einzudringen und sich etwas zusammenzubrauen.

Mit seinem Stock schwang er die angelehnte Tür auf und marschierte in ihr Inneres. Der Geruch nach Drachenbohnen stieg ihm in die Nase und ließ seinen Magen knurren, der sich die letzten Wochen primär mit Apfelnüssen und Sternkraut hatte begnügen müssen. Leander ignorierte den Eindringling in seinem Heim, so wie der Eindringling ihn ignorierte. So stapfte der Seher am Kamin vorbei, vor welchem er ein unheilvolles Blubbern vernahm, legte seinen Mantel ab, schlüpfte aus dem darunter liegenden varatanischen Brustpanzer und warf sich in seinen Schaukelstuhl.

„Nicht dorthin!“, ertönte Rekas warnende Stimme, „Maro genießt das Schlafplätzchen.“

Leander korrigierte seine Flugbahn und schrammte bloß gegen die Lehne des Schaukelstuhls, woraufhin ein wütendes Zischen von Rekas Schlange ertönte.

„Ich hoffe, ich störe dich nicht, Leander“, sprach Reka fröhlich weiter, „Ich brauchte ein Plätzchen, um dieses Rezeptlein zu perfektionieren, und du besitzt nun mal einen der besten Zwergenstahlkessel, die dieses Land zu bieten hat. Als du nicht anzutreffen warst, dachte ich, ich riskiere es. Ertappt!“

Leander konnte sich ein Zucken seiner Mundwinkel nicht verkneifen, dass er das Schweigeduell gegen die Kräuterhexe gewonnen hatte. Und er wusste es besser, als einen Streit mit Reka über die Nutzung seines Kessels anzuzetteln. Die Hexe wäre doch mitten in der Nacht in seine kleine Hütte gestürmt und hätte ihn aus dem Bett geschleudert, wenn ihr der Kopf danach stünde, seinen Kessel zu benutzen. Immerhin hinterließ sie üblicherweise als Ausgleich dafür ein kleines Geschenk wie ein Tränkchen oder eine uralte Schriftrolle, die Leander dann beim Baum der Lieder entziffern lassen konnte. Sie wusste, an welchem Wissen Leander sich erfreute, und so tolerierte er ihre Gegenwart. Auch wenn sie für seinen Geschmack zu viel vor sich hin brabbelte.

„Ich bin auch hier!“, meldete sich eine helle Stimme zu Wort, die definitiv nicht zu Reka gehörte, „Ich bin Chada, Rekas Schülerin. Also, vermutlich wisst Ihr das ja bereits. Reka meinte, es wäre in Ordnung, wenn...“

Leander blendete den Rest von Chadas Geplapper aus und knirschte die Zähne zusammen. Chada, die Heldin von Andor, deren herannahende Pfeile das letzte gewesen waren, was Taroks glutrote Augen erblickt hatten. Chada, die Heldin von Andor, die einen Kundschafter der Krahder aus dem Süden lieber erschlagen hatte, statt ihn gefangenzunehmen und nach Informationen auszufragen. Chada, die Heldin von Andor, die selbstgerecht Hademars Kenntnisse über die Verbindung von Dunkler Magie und Dunkler Hexerei aus dieser Welt gelöscht hätte, wenn sein Geist nicht glücklicherweise ohne seinen Körper hätte überleben können.

Nein, Leander und Chada hatten wahrlich nicht dieselbe Weltanschauung. Auch wenn Leander natürlich froh darüber war, dass die Helden von Andor diese Lande vor großen Übeln verteidigten, konnte er sich einfach nicht mit ihrer stolzen Zerstörungswut anfreunden. Es war, als kümmerten sie sich nicht einmal darum, dass mit jedem Streich ihrer Schwerter und jedem Pfeil im Herzen eines Gegners bessere Handlungsoptionen verloren gingen, tiefgehendes Wissen verloren ging, die Chance auf eine bessere Zukunft schwand. Und Chada war die stolzeste Heldin von allen, mit ihrem ehrwürdigen Auftreten und ihrer unüberwindbaren Willenskraft. Leander schüttelte seinen Kopf.

„Bis morgen seid ihr drei verschwunden, in Ordnung? Und bis dahin möget ihr eure Worte reduzieren. Mir steht der Kopf gerade nicht nach einem Schwatz“, murmelte Leander müde und ließ sich in einen anderen Stuhl sinken.

„Wann tut er das je?“, lachte Reka und warf irgendetwas in den Topf, woraufhin das ‚Wumpf‘ einer semi-großen Stichflamme zu vernehmen war.

„Fackel mir nicht das Dach ab!“, konnte sich Leander nicht verkneifen, „Das Rietgras ist staubtrocken und dieser Unterschlupf hat mich einiges gekostet.“

Die Hexe kicherte und schlug mehrmals mit einem metallischen Gegenstand gegen den Kochtopf.

Leander versuchte, seine strapazierten Glieder etwas zu entspannen, doch konnte er in der Gegenwart der hektisch umherhantierenden Hexe und ihrer etwas leiser, aber nicht weniger hektisch umhereilenden Schülerin kaum zur Ruhe kommen. Innerlich brannte er schon darauf, den Edelstein in seiner Manteltasche zu untersuchen, doch konnte er das kaum, während die alte Reka in seinem Haus umherhuschte.

„Ein sonderbarer Stein ist es, den du da bei dir trägst“, unterbrach Reka seinen Gedankengang, „Willst du ihn mir mal zeigen? Ich könnte dir verraten, welche Formen der Schatten unter seiner Oberfläche nimmt. Das ist wohl eine Grundlage für eine erfolgreiche Kommunikation mit der darin eigesperrten Seele.“

Wovon wusste diese Hexe denn nicht Bescheid?! Leander winkte dankend ab und zog zum Beweis das Drachenherz hervor. Chada machte ein staunendes Geräusch, auch wenn sie unmöglich wissen konnte, was in diesem Stein vorging.

„Sei still und lausche“, sprach Leander zu Reka, „und du wirst das leiseste Summen vernehmen. Schon bald wird er den Stein soweit unter Kontrolle haben, dass er ihn vibrieren lassen und so zur Kundgebung verwenden kann. Augen sind nicht vonnöten, um Töne zu hören. Selbst, wenn er sich nicht differenziert auszudrücken vermag, könnte ich ihm zum Beispiel den Code des Seefahrers Morseus beibringen. Ich brauche dich nicht, um mit ihm zu sprechen.“

„Wie du meinst“, antwortete Reka und klatschte irgendetwas, das nach einem großen Kochlöffel klang, in den Kessel. Sie flüsterte irgendwelche Anweisungen zu Chada, welche auf eine bestimmte Art im Kessel zu rühren begann, während Reka ein bisschen weiter weg in etwas Raschelndem herumkramte.

Leanders Gedanken schweiften wieder zu Varatans Fluch. Heute würde es zu spät sein, aber am nächsten Morgen könnte er zum Baum der Lieder aufbrechen und einen der dortigen Priester nach dem Stammbaum des Seekönigs fragen. Bestimmt wusste jemand etwas darüber, wie viele Kinder und Kindeskinder Varkmar gehabt hatte.

„Eigentlich“, setzte Reka an und unterbrach Leanders Gedankengang, „Eigentlich gibt es noch einen Grund, warum ich hier bin.“

Leander hätte sich schon innerlich an die Stirn geschlagen, wäre da nicht Rekas urplötzlich ernste Stimme gewesen, die ihn aufhorchen ließ.

„Eine alte Frau mit zu viel Zeit kommt viel um in diesen Landen, und da vernimmt man natürlich die einen oder anderen Gerüchte.“

„Was für Gerüchte hast du vernommen, Reka?“

„Man munkelt, Meres marschiere wieder durch diese Gefilde“, stieß Reka hervor, „Du hast ein Talent dafür, mächtige Gestalten aufzuspüren, Leander. Falls du ihn sehen solltest, könntest du ihm ausrichten, dass ich ihn gerne sehen würde? Es gibt da etwas, was ich zu berichtigen habe.“

Leander legte seinen Kopf schief. Das war ja aufrichtig interessant. Meres der Hexer war seines Wissens nach Rekas Schüler gewesen, dessen Hexerei die Hexenkunst Rekas bei weitem übertreffen konnte, aber stets unvorhergesehene Folgen mit sich zog. Meres hatte sich vor mehr als einem Jahrdutzend mit seiner Lehrmeisterin zerstritten – Rekas Hütte war dabei draufgegangen – und war im Zorn außer Landes gezogen. Nach dem Tod des Drachen war er zurückgekehrt, hatte ein Dutzend Tulgori aus dem bislang unbekannten Land westlich des Fahlen Gebirges mitgebracht, aus Versehen die Statuen des Bruderkrieges zum Leben erweckt und sein Leben im Kampf gegen diese Kreationen aus Stein riskiert, vielleicht gar gelassen.

Leander war nie davon überzeugt gewesen, dass Meres damals in Cavern umgekommen war. Er konnte es dem Hexer nicht verübeln, die Helden von Andor damals verlassen zu haben, nachdem sie ihn keines Blickes gewürdigt hatten. Und Reka hatte recht: Wenn Meres nun zurückgekommen war und Reka und die Helden mied, so hatte Leander wahrscheinlich die beste Chance, ihn aufzuspüren. Leander konnte den Mann zwar nicht gut durchschauen, aber er wusste, dass es einen Weg geben musste, ihn nach Narkon zu schicken. Würde ein so mächtiger Hexer wie Meres es schaffen, eine solche Kampfstärke gegen Varatans Fluch aufzubringen, dass auch der Wille aller Nachkommen Varatans gemeinsam nicht ausreichen würde, um ihn weiter wirken zu lassen?

Reka sprach hastig weiter: „Nein, tu... vergiss einfach, was ich gesagt habe. Es freut mich, dass du wieder im Lande bist, Leander.“

Die Kräuterhexe in Verlegenheit? Das war auch eine Seltenheit. Leander versuchte vergeblich mit seinem müden Geist, sich einen Reim darauf zu machen. Reka stapfte zur Tür hinaus, woraufhin Maro, die treue Seele, von Leanders Schaukelstuhl plumpste und der Hexe zischend hinterherschlängelte.

Chada wandte sich Leander zu, entschuldigte sich wortreich dafür, in seine Hütte eingedrungen zu sein, versprach, dass sie ihm dafür einen Gefallen schuldig wäre, und schlich dann leise zur Tür hinaus, einen netten Abschiedsgruß auf den Lippen.

Leander blieb noch einige Momente in seinem Sessel, um sicherzugehen, dass die Hexe nicht einfach kurz Wasser lassen war. Als sie und ihre Begleiterinnen nach dem Viertel einer Stunde noch nicht zurückgekehrt waren, übermannte ihn allerdings seine Neugierde und er erhob sich, um vor seinen Kamin zu schlurfen und nach dem Kessel zu tasten.

Der Dampf über dem Kessel roch lecker und unbedrohlich. Leander traute sich, einen Löffel in die Flüssigkeit zu tauchen und daran zu nippen. Seine Geschmacksnerven frohlockten, als er in Rekas Gesüpp einen der besten Eintöpfe erkannte, die er in seinem Leben je gekostet hatte.

Hatte sie ihm die ganze Zeit einfach nur ein Gericht gekocht? Ein Geschenk Rekas dafür, dass er Meres übermittelte, dass Reka ihn sehen wollte?

Versteh einer diese Hexen.

Nichtsdestotrotz war Leander dankbar für ihre Gabe.\bigskip







Wie geplant reiste Leander am nächsten Tag an den Baum der Lieder. Hohepriesterin Gända hatte frei und war so nett, ihn ins Archiv zu führen.

„Sagt noch einmal, nach wem Ihr sucht.“

„Varatan, der Seekönig. Ich nehme doch stark an, dass Ihr hier Aufzeichnungen über seine Familie lagert?“

„Selbstverständlich“, erwiderte Gända und Leander hörte, wie sich ihre Schritte hurtig entfernten.

„Warum interessiert Ihr euch für den alten Meereskönig, wenn ich fragen darf?“, erklang ihre durchdringende Stimme von der anderen Seite eines Regals.

„Ihr Jungspunde und Eure Neugier“, rief Leander aus. Er mochte es, einen mürrischen alten Gelehrten zu spielen und er wusste, dass diese Bezeichnung Gända ein bisschen Perspektive verleihen konnte. Gända war eine der altehrwürdigsten Bewahrer des Ordens und sehr stolz darauf, doch konnte sie nicht leugnen, dass Leander den Baum der Lieder schon aufgesucht hatte, als sie noch nichts als eine junge Novizin gewesen war.

„Ich habe kürzlich in einem Manuskript von ihm gehört“, flunkerte Leander das Blaue vom Himmel herunter, „Es gab ja schon viele Meereskönige von Varatan, aber da sein Reich unter seiner Herrschaft neu aufblühte und gar nach ihm ‚Varatanien‘ genannt wurde, lässt sich sein außergewöhnlicher Einfluss auf die Kultur der Nebelinseln nicht verweigern. Halb Werftheim will doch mit ihm verwandt sein! Und doch ist mir nur von einem einzigen Kind bekannt, das er in die Welt gesetzt haben soll. So etwas weckt doch das Interesse eines jeden wahren Gelehrten, wie wir es beide sind.“

„Gefunden!“, ertönte Gändas kühle Stimme. Ihre leisen Schritte auf dem Holzboden inmitten des Baums der Lieder kamen wieder näher.

„Varkmar, der erste Zauberer des Feuers, war der erste und einzige Sohn des Varatan.“

So weit, so gut, ganz wie erwartet, dachte Leander. Ein einziger Sohn, der inzwischen wahrscheinlich ohnehin schon tot war.

„Fahrt fort, mein Kind“, sprach er angespannt.

„Nun, viel ist über Varkmar nicht verzeichnet“, murmelte Gända beschämt, „Es gab stark widersprüchliche Schriften über sein Leben, aus den stark gefärbten Perspektiven der großen Zaubererorden Hadrias. Ich muss Euch wohl nicht berichten, wie konfliktreich es damals in Hadria zuging.“

„Irgendwelche neutralen Perspektiven?“, fragte Leander.

Er hörte es rascheln.

„Es gäbe hier einen Bericht, den die Heldin Eara aus dem Hohen Norden einst niedergeschrieben hatte. Das war zur Zeit, als sie einige Zeit lang hier in diesen Archiven hauste und unsere Aufzeichnungen zu den verschiedensten Zaubersprüchen studierte“, hängte Gända nicht ohne Stolz in der Stimme an.

„Eine Zauberin des Turms, und dann auch noch lange Zeit, nachdem es geschah?“, fragte Leander. Es war verblüffend, dass Gända auf eine solche Quelle überhaupt Wert legte.

Ein wenig spitzer als vorher fuhr Gända fort: „Nun, Eara ging sehr differenziert auf... jedenfalls berichtet sie von Varkmars Familie. Er hatte eine Tochter und einen Sohn.“

Leander horchte auf. Bei zwei Enkeln lag es auch noch im Bereich des Möglichen, beide Blutlinien aufzuspüren und auszulöschen.

„Und ihre eigenen Familien?“

„Nun, die Tochter...“, raschelte Gända erneut in den Pergamenten herum, „Die Tochter, Nika, verstarb jung, als Varkur das Siegel zum Eisernen Turm sprengte und die magischen Waffen zu erringen versuchte.“

Leander kam schwach eine Erinnerung auf, dass er diese Geschichte vermutlich schon einmal gehört hatte. Eine menschliche Schwäche war es, auf die er nicht stolz war, aber seitdem er Hadria verlassen hatte, hatte er sich nicht mehr groß um diese abgeschiedenste Insel des Hadrischen Meeres gekümmert. Er verband einige gute Erinnerungen mit dem Ort, aber noch mehr schlechte. Das hatte das traumatische Miterleben eines massenmörderischen Mobs so an sich. Leander schauderte beim Gedanken daran, wie Orweyn einst durch Hadria gewütet war.

„Und der Sohn?“, fragte er.

Gända stockte, und meinte dann: „O Leander... der Sohn Varkmars ist doch Varkur, der Dunkle Magier. Wusstet Ihr das nicht?“

Leander blieb stumm. Dann schlug er sich gegen seine Stirn. Wie hatte er so blind sein können?! Flammensohn, Flammenmeer, Flammentod! Es ergab solchen Sinn, Varkur passte so natürlich in die Abfolge.

Leander hatte von Varkur bislang nur als einen sich der Dunklen Magie verschrieben habenden Zauberer aus Hadria gedacht, der in der weiten Welt seine Chance auf Herrschaft und Rache suchte. Nun sah er, wie falsch er gelegen hatte. Varkur war weit mehr als das. Varkur trug das Erbe Varatans in sich, und stärkte den Fluch über Narkon mit seinem übermächtigen Willen.

Leander murmelte einige Dankesworte an Gända und verabschiedete sich von ihr. Seine Gedanken waren nicht mehr bei der Sache, sondern kreisten nur noch um einen: Varkur, den Enkel Varatans. Schon seit jungen Jahren aus Hadria verbannt. Höchstwahrscheinlich kinderlos. Und ungemein gefährlich.

Varkur musste sterben, damit Callem befreit werden konnte.

Und Leander würde dies irgendwie einrichten müssen.








\newpage
\section{Alte und neue Pläne}


Während Hademar im Laufe der nächsten Wochen langsam lernte, mit seinem kristallenen Gefängnis umzugehen und das Drachenherz zum Ruckeln zu bringen, verbrachte Leander den Großteil seiner Zeit am Baum der Lieder. Die Hoffnung, Callem zu befreien, hatte neuen Elan in sein vorher teilweise zielloses Leben gebracht und ließ ihn nun Gända, Tion und alle anderen Bewahrer, die sich seiner annehmen wollten, nach Texten über Varkur ausfragen.

Varkur war Leander vorher nie sonderlich interessant erschienen – ein Magier aus dem Hohen Norden, der hier und da auftauchte und vergeblich versuchte, seine Macht zu mehren. Doch nun, da Leander von Varkurs Verbindung zu Varatan wusste, war jeder noch so kleine Hinweis auf seinen möglichen Aufenthaltsort wertvoll.

Seinen letzten großen Auftritt hatte der Dunkle Magier gehabt, als er in der Ära des Sternenschilds den Helden den Sternenschild vor der Nase wegzuschnappen versucht hatte. Dieser letzte Angriff war wohl ein verzweifelter Plan des Hadriers gewesen, sein Schicksal zum Guten zu wenden, nachdem die Helden Tarok aus dieser Welt befördert hatten. Einer ziemlich ausführlichen Aufzeichnung über die erste Eroberung der Rietburg zufolge war Varkur überhaupt erst nach Andor gekommen, um den Drachen Tarok zu zähmen und gen Hadria zu reiten. So hatte er Rache an den Zauberern nehmen wollen, die ihn einst verstoßen hatten. Doch sein Hadrisches Stundenglas war zerborsten und Varkur hatte ein neues Ziel für seinen Hass gefunden: Brandur und die Helden von Andor. Hier und dort hatte er versucht, Taroks Kreaturen umzulenken und für seine Pläne zu missbrauchen, aber nie waren diese Pläne von Erflog gekrönt gewesen. Und dann waren Brandur und Tarok gestorben, und Varkur hatte plötzlich zwei große Ziele weniger gehabt.

Nun, nachdem die Helden ihm den Sternenschild entrissen hatten, hatte er wohl definitiv einen neuen Plan finden müssen, mit dem er seine gewünschte Rache an den Helden und an Hadria vollziehen könnte. Aber seltsamerweise hatte er seit damals keine weiteren offen Attacken mehr durchgeführt.

Mit gefurchter Stirn ging Leander aber und abermals die gehörten Textstellen in seinem Kopf durch. Varkur befand sich höchstwahrscheinlich nicht mehr in Andor, sonst hätten die Helden ihn bereits aufgespürt. Aber er würde doch auch nicht einfach irgendwo ins Gebirge geflohen und dort gestorben sein, oder? Er musste ein Versteck besitzen und an seinen nächsten Plänen feilen, und was seine Pläne auch waren, es mussten große sein. Nun, wo Tarok nicht mehr war, wagte sich Varkur vielleicht an den Urtroll? Könnte er diesen etwa ohne zusätzliche Zeit eines hadrischen Stundenglases zu bändigen lernen? Oder hatte er sich analog zu Meres ins ferne Tulgor zurückgezogen, um dort neue Technologien zu meistern?

Leander griff gedankenverloren nach dem Drachenherz und schüttelte den Edelstein, woraufhin ein leises Klingen ertönte. Hademar reagierte!

„Was meinst du, wo sich dein Varkur aufhält?“, fragte Leander den Stein leise. Hademars Singsang schwankte unkontrolliert und brach dann ab. Noch war er nicht soweit, sprechen zu können. Aber das war in Ordnung. Gut Ding will bekanntlich Weile haben. Früher oder später würde er rudimentär mit der Außenwelt kommunizieren können. Und Leander würde da sein, um ihm zuzuhören.

Hademars Geist war tatsächlich eine von Leanders letzten Hoffnungen darauf, Varkur aufzuspüren. Schließlich waren der aus Krahd geflohene Nekromant und der aus Hadria verbannte Dunkle Magier lange Zeit Komplizen gewesen. Die genaue Natur ihrer Beziehung blieb Leander fürs Erste verborgen, da kaum mehr als ein einzelner von den Helden abgefangener Brief Hademars an Varkur von ihrer Bekanntschaft zeugte. Wenn diesem uneingeschränkt Glauben zu schenken war, so hatte Hademar nach seiner Flucht aus der Gefangenschaft der Krahder dem Land Andor, das von seinem ihn zurückgelassen habenden Bruder Brandur regiert wurde, den Rücken zugekehrt und war nach Hadria aufgebrochen. Die Krahder hatten sein Talent für die Dunkle Hexerei geweckt gehabt und nun hatte es Hademar danach gedürstet, mehr zu erlernen.

Doch die Zauberer des Turms hatten Hademar abgewiesen. Kein Wunder, diese elenden Narren fürchteten sich doch vor allem, das potentiell gefährlich sein konnte, also nebst ihren eigenen Schatten auch vor so gut wie allen frischen Ideen. Im aus der Feste von Yra verstoßenen Varkur hatte Hademar allerdings einen Seelenverwandten gefunden. Die beiden hatten viel Zeit in Nordgard verbracht und Rachefeldzüge gegen die Zauberer der Feste von Yra geplant. Varkur hatte Hademar Zugang zur Bibliothek der Akademie von Hadria verschafft und Hademar Varkur vom Drachen Tarok tief unter dem Grauen Gebirge erzählt. So waren sie gemeinsam ausgezogen. Während Varkur den offenen Kampf mit Andor gesucht hatte und prompt unterlegen war, hatte sich der Pfad des vorsichtigeren Hademars schon früher von Varkurs getrennt. Der Nekromant hatte sich zunächst in der Winterburg eingerichtet, wo er die umliegenden Agrenvölker versklavt und langsam Drachenmagie aus den Zwergentürmen im Grauen Gebirge gezogen hatte. Letzten Endes war Hademar erheblich stärker gewesen als Varkur, aber selbst das hatte ihn nicht vor dem Zorn der Helden retten können.

Und nun lag Hademars Geist in Leanders Hand, wusste vielleicht, wo Varkur sich aufhielt, konnte es Leander aber nicht verraten. Noch nicht. Doch Leander war geduldig. Er konnte warten. Und Varkur würde ihm nicht davonlaufen.\bigskip







Ironischerweise spürte Leander nie Varkur auf. Stattdessen spürte Varkur Leander zuerst auf.

Das erste, was Leander spürte, war ein leiser Schauer, der seinen Arm entlangglitt. Er richtete sich in seinem Bett auf und hoffte, dass es nicht schon wieder Shan wäre, die ihn in seinem Schlaf störte.

Es war nicht Shan. Soviel war urplötzlich klar. Leander spürte ihn herannahen. Jedes Haar an seinem Körper stellte sich auf und eine Schwere überkam seinen Körper von allen Seiten zugleich. Dann begann die Pein. Die gesamte Länge von Leanders Arm verkrampfte sich und flammende Schmerzen strahlten bis in seinen Kopf. Leander hatte genug von Varkurs Angriffen gehört, dass er dieses Phänomen als den schmerzinduzierenden schwarzen Nebel erkennen konnte, mit dem sich der Dunkle Magier nur allzu gerne umgab.

Mit zusammengebissenen Zähnen versuchte Leander, den Schmerz zu ignorieren und seinen Gegner zu lokalisieren. Er drehte seinen Kopf in die Richtung, aus der er Schritte zu hören glaubte, und stolperte dann in entgegengesetzter Richtung aus seinem Bett.

Für einen kurzen Augenblick ließen die Schmerzen nach, dann war der Nebel des Magiers wieder um ihn herum, noch stärker als vorher. Der Mantel der Agonie stülpte sich über Leander, riss den dürren Seher in die Höhe und schleuderte ihn zurück auf seine Schlafstelle. Leander hörte etwas knirschen. Hoffentlich war das nur das Bett gewesen und nicht sein Rücken.

Dann war da urplötzlich ein Gestank nach Blut und faulem Fleisch, den Leander sonst nur von Kreaturen kannte. Er brauchte einen Augenblick, um zu erkennen, dass das Varkurs rasselnder Atem war. Der Dunkle Magier befand sich direkt vor ihm!

„Wo ist er?“, sprach eine zischende Stimme leise und langsam, „Wo ist Hademar?“

Leander hätte beinahe aufgelacht. Wenn er gewusst hätte, dass Varkur sich noch um seinen ehemaligen Komplizen sorgte, hätte er Hademars Kristall noch viel prominenter als Lockmittel eingesetzt. Nur war dies eine gänzlich unerwünschte Konfrontation, denn auch wenn er Varkur hatte aufspüren wollen, war Leander im Moment planlos und hoffnungslos davon entfernt, Varkur zu töten. Seine Priorität musste nun darin liegen, diese Situation mit seinem Leben zu verlassen. So griff Leander ächzend in seine Manteltasche und präsentierte Varkur gefügig das Drachenherz, den weißen Runenstein mit Hademars Seele darin.

„Töte mich nicht, o Sohn des Varkmar“, stieß Leander hervor, „Ich bin der Einzige, der Hademars Geist wieder aus diesem Stein befreien kann.“

„Du unterschätzt mich“, hauchte die kalte Stimme Varkurs. Schuppige Klauen klaubten Leander das Drachenherz aus klammen Fingern. Dann verdichtete sich der schmerzvolle Griff von Varkurs Nebels um Leander Kehle. Dieser ächzte auf und schnappte nach Luft, versuchte vergeblich, weitere Worte hervorzubringen, doch stattdessen floss nur noch mehr des betäubenden dunklen Nebels in seinen Hals hinein und stoppte seine Luftzufuhr.

Leander wusste, dass er nur wenige Augenblicke bei Bewusstsein bleiben konnte. Er musste handeln, und zwar schnell. Im direkten Kampf gegen den Dunklen Magier hatte er keine Chance, aber so...

Mit seiner letzten Kraft hob Leander seine Hand und stach mit spitzen Fingern dorthin, wo er Varkurs Gesicht vermutete. Er schrammte gegen Schuppen und Horn, doch dann traf er auf etwas Glibbriges, das Leander als Varkurs Augen identifizieren konnte.

Der Dunkle Magier heulte auf und schleuderte Leander von sich. Wohl wissend, dass dies seine letzten Sekunden in dieser Welt sein könnten, ignorierte Leander seinen protestierenden Körper, hob sich in die Höhe und stolperte zur Tür hinaus in den nächtlichen Wachsamen Wald hinein.

Schon hatte sich Varkur von der Überraschung erholt und brüllte ihm etwas dämonisch Klingendes hinterher. Leander indes erhob seine eigene Stimme und schrie:

„Hilfe! Hilfeeeeeee!“

Der Wachsame Wald blieb stumm, bis auf eine Eule, welche weit entfernt antwortete. Leander wägte ab, ob er erneut schreien sollte, entschied sich dann dagegen, weil dies Varkur auf seine Position aufmerksam machen könnte, und rannte vom ihm bekannten Waldpfad weg.

Blind stolperte der Seher durch das Unterholz, schrammte hier und dort gegen Äste und Wurzeln, die ihm ins Gesicht schlugen und seine Füße beinahe zum Stolpern brachten. Beinahe.

Die Hände ausgestreckt, versuchte Leander, sich zumindest ein klein wenig vor Hindernissen zu bewahren. So kam es, dass Leander immerhin nicht kopfsvoran mit einem Baum kollidierte, sondern den Aufprall etwas abfedern konnte. Trotzdem schleuderte ihn dieses unglückliche Treffen zu Boden, wo er hustend und prustend nach Luft rang, während die Schmerzen aller Schnitte und blauen Flecken, die der Wald ihm verpasst hatte, auf ihn einprasselten und nur schwer zu ignorieren waren.

Wie weit hatte er sich von seiner Hütte entfernt? War Varkur ihm auf den Fersen?

Seine Brust hob und senkte sich schnell, während Leander seinen Atem kontrolliert beruhigte und sich so flach wie möglich an den Baumstamm kauerte, gegen den er gerannt war.

Vielleicht war er Varkur entkommen. Vielleicht war dieser so sehr auf Hademar fokussiert, dass ihn Leander kaum kümmerte, oder zumindest so lange nicht kümmerte, dass Leander hatte entkommen können.

Leander Herz sank in sein niederes Nachtgewand, als sein Körper von einem erneuten Schauer erfasst wurde und etwas Schmerzvolles, Nebliges sich um seine Hüfte schnallte. Mit einem Ruck wurde er aus dem Unterholz gerissen und strampelte vergeblich. Seine Füße streiften kaum mehr die oberen Enden der Sträuche.

„Du dachtest doch nicht wirklich, dass es so leicht wäre, einem Dunklen Magier zu entkommen?!“, lachte Varkur finster. Seine Stimme klang ein wenig weiter weg als vorhin. Er hatte gelernt, Leanders Finger seinen Augen nicht zu nahe kommen zu lassen.

„Ich... ich kannte deinen Vater“, stieß Leander verzweifelt hervor, „Und deinen Großvater! Ich könnte dir von ihren Geheimnissen verraten, dir Macht liefern. Hademar weiß, wie man den Tod umgehen kann. Wünschst du dir nicht, deine Schwe...“

Für einen kurzen Blick hatte sich der Griff um Leanders Körper gelockert, doch nun wurde er noch fester als vorhin. Varkur brüllte hasserfüllt: „Ich habe keine Väter! Und ich wünsche mir keine Händel mit dem Tod! Ich bin eine Kreatur der Finsternis! Alles Menschliche in mir wurde schon vor Urzeiten ausgelöscht!“

„Mit Verlaub, das ist eine Unwahrheit“, erklang eine tonlose aber nicht unfreundliche Stimme. Leander kannte sie nicht. Varkur offenbar auch nicht, denn der Dunkle Magier fuhr herum und knurrte:

„Wer da?!“

Der Nebel des Schmerzes um Leander löste sich plötzlich in Luft auf, woraufhin der Seher unsanft auf den Boden knallte. Sein Schädel brummte, doch stemmte er sich ungeachtet dessen in die Höhe und stolperte so rasch wie nur möglich von Varkur weg.

„Dunkle Magie ist in diesen Landen nicht gern gesehen. Und nächtliche Überfälle erst recht nicht. Ich muss dich bitten, dich unverzüglich zu ergeben, Unruhestifter“, erklang dieselbe monotone Stimme erneut. Kein Bewahrer wäre so dreist, sich Varkur direkt im Kampf zu stellen.

Varkur lachte auf und Leander hörte einen gellenden Aufschrei. Vermutlich hatte Varkur seinen Nebel auf den Neuankömmling fokussiert. Leander dankte seinem scheinbar todesverliebten Helfer innerlich und zog sich vorsichtig weiter aus dem Kampf zurück. Da erklang das ‚Wumpf‘ einer Stichflamme und der Schrei des Neuankömmlings versiegte. Stattdessen vernahm er, wie Varkur auffluchte.

Etwas knisterte und den Gestank verbrannten Fleischs drang an Leanders strapazierte Nase. Dann ein weiteres ‚Wumpf‘, das Rascheln von Stoff und... waren das Katzen? Ja, das war eindeutig das Miauen dutzender Katzen, direkt von der Stelle, wo er den Kampf zwischen Varkur und dem offensichtlich magiebegabten Neuankömmling vermutete. Leander kratzte sich am Kopf.

Das Miauen ließ nicht nach, dafür das Geräusch magischer Attacken, welches urplötzlich versiegte. Leander duckte sich auf den Boden und versuchte, so unsichtbar wie möglich zu bleiben. Leise Schritte näherten sich ihm ungeachtet seines Versteckversuchs. Einerseits verriet ihm das, dass er trotz seiner Mühen von weitem sichtbar gewesen war. Andererseits wusste er, dass Varkurs Schritte nicht so klangen. So vermutete Leander, dass es der mysteriöse Neuankömmling war, der sich auf ihn zubewegte. Und er hatte schon eine Vermutung, um wen es sich handeln könnte.\bigskip







„Der Magier hat sich verzogen. Ihr seid in Sicherheit“, erklang die tonlose Stimme des Fremden. Gedankenverloren fügte er hinzu: „Eigentlich wollte ich ihn in seinem eigenen Nebel eingangen. In der Steppe Tulgors fand ich eine Blume, deren Saft Wasser- und auch Säuretropfen blitzschnell einfrieren können. Ich hatte keine Ahnung, warum auf einmal all diese Katzen aufgetaucht sind. Eine faszinierende Reaktion...“

Seine Stimme brach ab, als sein Fokus sich auf etwas anderes richtete, sein Mantel und urplötzlich ein lautes Schnurren vom Boden ertönte. Die Pause erlaubte Leander, sich hoheitsvoll aufzurichten und zu vermuten: „Seid Ihr Meres, der Hexer von Andor?“

Sein Gegenüber klang das kleinste Häuchlein von überrascht, als es antwortete: „Den Hexer von Andor, so kennt man mich hier noch? Ja, ich bin Meres, oder zumindest war ich das einst. Ich bin mir nicht mehr sicher, wer ich bin. Meres. Haamun. Ein Schüler. Ein Hexer. Ein Reisender. Ein Unruhestifter. Ein Beschützer.“

„Das klingt alles sehr nobel“, bekräftigte Leander abgelenkt, während er sich sein malträtiertes Nachthemd glattstrich, „Ich bin Leander der Weise, und ich danke dir außerordentlich für dein tapferes Eingreifen. Ohne dich gäbe es mich wohl nicht mehr. Lass mich dich zum Dank in meine Hütte bitten und dir eine Gegenleistung überreichen.“

Meres blieb eine Zeit lang still und murmelte dann: „Die Helden von Andor verlangen keine Gegenleistung für ihre Dienste, oder?“

Leander blickte ihn überrascht an: „Meinst du etwa, du wärst... öhm, nein, höchstens einen deftigen Eintopf als Dank habe ich sie schon annehmen sehen. Und natürlich die königlichen Kopfgelder für überwundene Kreaturen.“

Unweigerlich schweiften Leanders Gedanken zurück zu demjenigen Tag, als er mit Hademar in der Manteltasche aus dem Grauen Gebirge zurückgekehrt war und Reka in seiner Hütte aufgefunden hatte. Er hatte ihr damals versprochen gehabt, Meres auszurichten, dass sie mit ihm reden wollte. Vermutlich hatte die Hexe im Sinn, ihren Zwist mit ihrem ehemaligen Schüler beizulegen, und wer konnte es ihr verübeln?

Es stand zu vermuten, dass Reka früher oder später davon erfahren würde, wenn Leander Meres Rekas Wunsch verschweigen würde. Sein Versprechen zu brechen, würde also in einem Verlust des wenigen bisschen an Vertrauens resultieren, das die Hexe noch in ihn steckte. War ihm das die kleine Chance wert, dass Meres Varatans Fluch überwältigen konnte?

Leander überdachte die Sachlage kurz und kam zum Schluss, dass ein Mittelweg das optimale Vorgehen wäre. Er würde Meres zunächst zu Reka schicken. Aber das bedeutete nicht, dass er ihn vorher nicht auf den richtigen Pfad locken könnte.

„Hör mal, Meres, Reisender und Beschützer. Nach was suchst du hier in diesem Lande?“, setzte er an.

Stille, bis auf das anhaltende Schnurren der Katzen um Meres herums. Dann...

„Ich weiß es nicht. Ich nehme an, ich suche meinen Platz auf dieser Welt.“

„Sehr poetisch, o Hexer. Hast du diesen Platz bereits gefunden?“

„Ich habe vieles gefunden, und vieles hinter mir gelassen“, kam die monotone Antwort wie aus der Achter- oder Arcuballiste geschossen, „Aber noch an keinem Ort fühlte ich mich wirklich zuhause. Nicht viele verstehen mich. Vielleicht ist das mein Los.“

Sehr schön, dachte Leander innerlich und fuhr fort.

„Für einen jeden gibt es eine Bestimmung. Die richtigen Menschen für einen zu finden, kann den Unterschied zwischen einem leeren und einem erfüllten Leben machen.“

Leander verdrängte die Gedanken an die brüderliche Verbindung, die er und Callem geteilt hatten. Nicht einmal von Reka hatte er sich je so verstanden gespürt. Aber jetzt musste er sich auf Meres konzentrieren:

„Vielleicht könnte ich dir helfen, herauszufinden, wo du dich auf die Suche machen musst. Ich glaube zu spüren, wo sich eine gute Gemeinschaft für dich aufhalten könnte. Ich bin nämlich ein Seher.“

„Na sieh einer an, ein Seher?“, erwiderte Meres höflich, „Ich halte es für unwahrscheinlich, dass du mir zu zeigen vermagst, was nicht einmal ich selbst herauszufinden vermochte.“

„Na komm, was hast du schon zu verlieren?“

Und Leander packte Meres bei den Schultern und führte ihn zu seiner Hütte zurück.\bigskip







Varkur ließ sich langsam zu Boden sinken. Der schwarze Rauch der Dunklen Magie setzte sich um ihn herum und färbte den Schnee des Grauen Gebirges in einem dunklen Ascheton.

Er musste nießen. Schnee und Rauch stoben gleichermaßen auf und verborgen den Magier von potentiellen neugierigen Blicken.

Katzen. Ausgerechnet Katzen hatte dieser unvorsichtige Hexer herbeirufen müssen! Wie konnte das sein, dass Varkur eine solche Macht besaß, und seine schon beinahe vollständig echsenhafte Nase noch immer so stark auf diese elenden Biester reagierte?!

Erneut erschütterte ein Niesreiz Varkurs Körper und sein magischer Nebel trug die Vibrationen in seine Umgebung. Er würde zurückgehen müssen und sich um diesen halbgaren Hexer kümmern, der glaubte, mit Kräutern und Tränklein an seine Künste ankommen zu können. Er würde ihm lehren, was es hieß, sich auf ein Duell einzulassen mit Varkur, dem größten Dunklen Magier seit Orweyns Zeiten! Aber im Moment war das nicht seine Priorität.

Hatschi!

Wütend vor sich hin murmelnd riss Varkur einige Sternkraut-Büschel aus dem Untergrund und zerquetschte sie in seiner Faust, bis ein grünlicher Saft aus dieser heraustropfte. Varkur schmierte sich die nasse Masse in die tränenden Augen und zuckte zusammen, als sein Körper ihm wieder schmerzlich daran erinnern ließ, dass der Seher ihn kürzlich in ebenjene gereizten Augen gestochen hatte. Immerhin konnte der Sternkraut-Saft den Juckreiz etwas lindern.

Varkur hasste es, einen Körper zu haben. Ständig schmerzte etwas oder ging zu Bruch. Ein freier Geist, so wie Hademar es ihm vorgeschwärmt hatte, das wollte er sein. Aber noch besaß er nicht die nötigen Kenntnisse dazu, körperlos fortzubestehen. Noch musste er mit dieser Hülle klarkommen.

Varkur schloss seine schmerzenden Fenster zur Seele und tauchte erleichtert in eine angenehme Dunkelheit ein. Ohne groß nachzudenken, dirigierte er den schwarzen Rauch der Dunklen Magie, ließ ihn zerfasern, sanft die Umgebung ertasten, ihm ein Bild übermitteln.

Nicht weit von ihm entfernt krochen zwei Gorlots aus einer Höhle, nichts von seiner Gegenwart ahnend. Diese gelenkigen Gestalten waren weitaus geschickter als gewöhnliche Gors (kein Wunder, sie hatten ja Hände) und waren dazu noch gescheiter. Aber sie hatten weder die natürliche Rüstung von Gors noch das taktische Geschick von Skralen, von der Geschwindigkeit von Warkreaturen oder der Stärke von Trollen ganz zu schweigen. Sie waren schlicht und ergreifend zu gewöhnlich, als dass sie Varkur in seinem Heer gewinnbringend einsetzen könnte, und Varkur konnte nur so und so viele Kreaturenanführer auf einmal dirigieren. Aber wenn gerade nichts Besseres zur Verfügung stand...

Varkur drang in den Geist des einen Gorlots ein und dieser wurde von einem einzigen Gedanken erfasst.

Der Gorlot erhob den tarenähnlichen Kopf mit dem übergroßen Unterkiefer in den Himmel und schnüffelte. Dann jaulte er auf und rannte auf Varkur zu. Sein Kumpane grunzte verwirrt und setzte ihm dann nach. In Windeseile waren die beiden Kreaturen bei Varkur angekommen, die eine willig wartend, die andere ein rostiges Messer zückend, ein Überbleibsel aus dem Unterirdischen Krieg, das der Gorlot sich angeeignet hatte.

Der dichte schwarze Nebel der Dunklen Magie machte kurzen Prozess mit dem zweiten Gorlot, während Varkur das Drachenherz hervorholte, den einst weißen, nun blutroten Edelstein mit beiden Händen fest umschloss und versuchte, die verworrenen Fäden der Drachenmagie zu durchdringen, die Hademars Seele im Kristall einsperrten. Als ihm das wie erwartet nicht gelang, versuchte Varkur stattdessen, das verworrene Geflecht der magischen Ströme aufzunehmen und mithilfe Dunkler Magie in den schwachen Geist des Gorlots zu projizieren.

Der Gorlot blieb weiterhin von Varkurs Befehl erfüllt stehen und kratzte sich gedankenverloren am Hinterbein. Mehr geschah nicht. Varkur fluchte auf und konzentrierte sich noch stärker. Schweiß tropfte aus dem menschlichen Teil seiner Haut und perlte an seinen Schuppen ab. Es kam ihm wie eine Ewigkeit vor, bis der Gorlot endlich aufjaulte und auf seinen ungeschützten Hintern plumpste, während das projizierte Geflecht des Geists im Drachenherzens seine Gedanken überschrieb und unterdrückte.

Varkur öffnete seine gereizten Augen und erblickte durch einen Tränenschleier hindurch den Gorlot, wie er sich überrascht umsah und dann sorgsam aufrichtete. Ein krächzendes Stöhnen drang aus seiner tierischen Kehle, als Hademar seine neuen Stimmbänder ausprobierte. Dann öffnete der Gorlot seinen unmenschlich großen Kiefer, entblößte mehrere Reihen spitzer Zähne und sprach:

„Varkur, du alter Sack! Wie hast du nur so lange gebraucht, um mich zu finden?!“

Nicht einmal zu einer Erwiderung ansetzen konnte Varkur, ehe die Verbindung abbrach, Hademars Geist wieder hinter die magischen Mauern des Drachenherzens gezogen wurde und der Gorlot zusammenklappte wie eine Puppe, der man die Fäden durchgeschnitten hatte. Die Kreatur richtete sich zuckend auf, sah so erschrocken drein, wie man mit einem solchen Gesicht erschrocken dreinschauen konnte, und sprang dann keuchend und stolpernd davon.

Varkur richtete seinen Blick auf das Drachenherz in seiner Hand. Es hatte ihn schon solche Unmengen an Energie gekostet, Hademars Geist einige Lidschläge lang den Gorlot steuern zu lassen. Wie in aller Welt wollte er dies über längere Zeit aufrecht erhalten?

Nein, eine langfristige Lösung war das wahrlich nicht.

Aber genug, um mit Hademar zu kommunizieren.

Genug, um herauszufinden, wie er ihn aus diesem Gefängnis befreien konnte.\bigskip







Leander war von gemischten Gefühlen erfüllt. Einerseits war er froh darüber, Meres von Narkon berichtet und den Hexer auf die verfluchte Insel angesetzt zu haben. Es fiel Leander uncharakteristisch schwer, Meres einzuschätzen, doch schätzte er seine Chancen, Meres‘ Interesse geweckt zu haben, nicht schlecht ein.

Andererseits hatte er, der stets stolz „Ich bin ein Seher“ verkündet hatte, es wieder einmal immer wieder versucht und nicht geschafft, auch nur den Ansatz einer Vision zu empfangen. Seit dem Gespräch mit der Stammesältesten hatte er tatsächlich keine Nachricht mehr von der Zukunft empfangen und das leise Gefühl von Unfähigkeit, gepaart mit der Angst, seine Fähigkeit verloren zu haben, setzte wieder ein.

Am Ende hatte er Meres einige Drachenknochen hingeworfen und ein bisschen Hokuspokus darum und um Narkon gesponnen, nur um danach auf Rekas Wunsch, ihn zu sehen, zurückzukommen. Das war mitnichten seine überzeugendste Schauspielrolle gewesen.

Ganz zu schweigen davon hatte er bei dieser ungewünschten Begegnung mit Varkur das Drachenherz und Hademars Seele verloren. Dem Ziel, Varkur zu töten, war er keinen Schritt näher gekommen und nun wusste er mit Sicherheit, dass der Dunkle Magier nicht aus purem Glück das Zeitliche gesegnet hatte. Und da Hademar verschwunden war, konnte er jetzt nicht einmal mehr den Nekromanten nach seinen Kenntnissen oder nach Varkurs Versteck aushorchen.

Leander hoffte, dass Varkur von Meres‘ unvorhersehbaren Hexereien derart abgeschreckt worden war, dass er nicht zurückkehrte, um Leanders Leben endgültig zu beenden. Aber Leander fürchtete sich nicht sonderlich vor einer Rückkehr Varkurs. Einerseits hatte Varkur keinen persönlichen Hass auf ihm, andererseits war Varkur im Moment wohl ohnehin damit beschäftigt, Hademars Geist aus dem Drachenherz lösen zu wollen. Nomion würde nicht glücklich sein, sobald es ihm gelang. Aber das war fürs Erste nur dessen Bier, und vielleicht eines Tages das der Helden von Andor. Leander wollte im Moment vor allem eines: In sein zusammengestürztes Bett fallen und sich heilen.

Ein bisschen krustige Heilsalbe hatte Leander sicher noch übrig. Er schleppte sich in seinen Trankraum, tastete die Regale nach der passenden eingekratzten Inschrift ab und zog dann das Säcklein hervor. Ächzend ließ er sich auf dem erdigen Boden seiner Hütte nieder, öffnete das Säcklein und befeuchtete das körnige Pulver mit etwas Spucke. Es schmeckte bitter. Alsbald tastete Leander seinen Körper nach offenen Wunden ab und schmierte etwas von der körnigen Paste auf alle gefundenen Schnitte und Schrammen.

Da floss eine unnatürliche Kälte in den Raum und ließ die Wundpaste einfrieren.

„Was ist denn jetzt schon wieder?!“, murmelte Leander missmutig und zog sich in die Höhe. Es war noch nicht einmal Morgen!

Wie als Antwort darauf schlängelte sich ein eiskalter Schattententakel um seine nackten Beine und zog ihn unsanft nach vorne. Seine Wunden öffneten sich wieder und der Tentakel zerfaserte sich in mehrere kleinere Ärmchen, welche sich ebenso unsanft um Leanders wunden Körper wanden.

„Shan“, gähnte Leander erschöpft, „Wir hatten doch bereits...“

Er verstummte erschrocken. Das Drachenherz, oder besser gesagt das darin gespeicherte Licht der Welt, das ihn bislang vor Shans Zorn bewahrt hatte, war nicht mehr hier. Varkur hatte es mitgenommen. Konnte es wahrlich sein, dass die seelenhungrige Shan Leander nachgeschlichen war, nur darauf wartend, dass ihn das Drachenherz verließ?

„Jetzt komm schon, die mächtige Seele des Nekromanten ist nicht hier“, sprach Leander unwirsch, in der Hoffnung, die Schattenhexe noch auf Varkurs Spur zu lenken, „Du spürst doch bestimmt, wohin der Dunkle Magier sie bringt. Wenn du dich beeilst, kannst du den Moment abpassen, wenn er die Seele freilässt.“

Äußerlich gab er sich gelassen, ja, leicht genervt. Innerlich musste er seine Panik unterdrücken, denn ohne einen Edelstein war er Shan schutzlos ausgeliefert.

Shan antwortete lächelnd: „Keine Sorge, o Seher. Ich werde dort sein, wenn der Dunkle Magier den Schutz vom Geist des Nekromanten entfernt. Aber zuerst wende ich mich dir zu. Es kommt nicht oft vor, dass mir jemand entkommt. Du hast mein Interesse geweckt und darum werde ich dich mir nun einverleiben. Der Schatten ist hungrig.“

Wie viel musste Leander denn diese Nacht noch erdauern?! Unauffällig tastete er nach seinem Holzstab auf dem Hüttenboden. Der Stab war kein Zauberstab, nur eine Replik, die Eindruck schindete und ihm bei der Fortbewegung half. Aber solange er als stumpfe Waffe taugte, musste er gar nicht magisch sein, um etwas anrichten zu können. Leander wusste, dass die Helden von Andor Shan bei der Befreiung der Rietburg schon einmal vertrieben hatten, und die meisten von ihnen besaßen kein ausgeprägtes magisches Talent. Insofern musste es möglich sein, Shan alleine durch körperlichen Schaden das Weite suchen zu lassen. Also musste es auch für Leander möglich sein, die Schattenhexe zu vertreiben.

„Suchst du das hier?“, erklang Shans leise Stimme, in der Leander nun einen Anflug von Häme zu erkennen glaubte. Ein helles Poltern erklang, als Shan Leanders Stab ergriff und an die gegenüberliegende Wand schleuderte.

Leander stieß sich an der Wand hoch. Sein Messer lag ebenfalls nutzlos weit entfernt. Meres war fortgezogen. Er stand mit dem Rücken zur Wand und Shans Nähe ließ ihn immer mehr frösteln. Schon spürte er die verfluchten Schattenkräfte Shans nach ihm greifen, nahm wahr, wie seine Kräfte schwanden und eine große Last sich auf seinen Geist legte.

Der Seher von Narkon stürzte verzweifelt nach vorne, hinein in die Kälte. Er griff mit klammen Fingen nach dem zerfasernden Körper Shans, versuchte, irgendetwas zu fassen zu kriegen, das er festhalten konnte, dem er Schaden zufügen konnte. Shans Präsenz fing ihn weich, beinahe sanft auf und stülpte sich dann von allen Seiten um ihn, schnürte ihm die Atemluft ab. Sternchen blitzten vor seinem inneren Auge auf und er sank auf seine Knie.

„Shan“, krächzte er auf, „Ich kann dir so viel nützlicher sein, wenn du mich fürs Erste am Leben lässt.“

Die Schattenhexe ließ nicht locker. Leander wusste nicht, was er weiter sagen sollte. Damals, im Kerker der Rietburg, hatte es gewirkt, als wäre Shan für solche Argumente empfänglich gewesen. Doch nun agierte sie nur ihrem Seelenhunger nach. Leander verstand den Drang der Dunkelheit, der auf der Schattenhexe lag, und er wüsste nicht, wie er sie umstimmen sollte, wenn sie nicht auf ihren Verstand hörte.

Wahrscheinlich würde dies sein Ende sein. Vielleicht hatte er deswegen nur noch so selten Visionen empfangen in der letzten Zeit.

Leanders letzter Gedanke galt Callem.\bigskip







Da öffnete sich sein sehendes Auge und Licht strömte auf ihn ein. Leander rang nach der frischen Luft, so illusorisch sie auch sein mochte, und erblickte ein kleines Mädchen mit einem lieben Lächeln auf den Lippen, wie es durch den freien Markt hüpfte. In seinen Händen trug es einen Korb mit frischen roten Äpfeln. Leander zog die satten Farben der Früchte ein.

Seine Mundwinkel zuckten unwillkürlich nach oben beim Anblick dieses unschuldigen Lächelns. Vergessen war der klamme Griff Shans, der ihn in der Realität gepackt hatte und ihm unzweifelhaft gerade das Leben aussog.

„Rosanna! Sanna! Komm, meine Sonne!“, erklang eine freundliche Stimme. Das kleine Mädchen blickte auf und rannte zu einer großgewachsenen Frau, welche sie hochhob und fröhlich im Kreis herumschwang, einmal, zweimal, dreimal.

Da flackerte das glückliche Bild und an der Stelle des kleinen Mädchens sah Leander eine von flirrenden Schatten umgebende Gestalt mit einer finsteren Fratze. Shan war das, und Shan öffnete ihren Mund. Alle Schatten der Welt traten daraus hervor, überfluteten den freien Markt und schwemmten die fröhliche Frau davon. Die Farben des Lichts ergrauten, die Sonne schien ohne Wärme. Nichts blieb mehr übrig außer Shan, die inmitten der Dunkelheit stand und verzweifelt aufschrie. Einen solchen Schrei der Verzweiflung hatte Leander erst kürzlich bei der Stammesältesten vernommen. Sein ganzer Körper erschauerte ob des Klanges.

Oder erschauerte er wegen der Kälte, die immer noch von Shans Griff ausging? Das innere Bild verblasste und hinterließ die übliche Absenz jeglicher optischer Sinnesreize. Er hatte eine weitere Vision gehabt. Aber keine Vision der Zukunft, sondern eine der Vergangenheit. Das war neu. Hatte irgendetwas an Shans einzigartiger Schattenmagie sein sehendes Auge verwirrt und in die Vergangenheit statt die Zukunft blicken lassen? Es war egal.

Rasch zählte Leander Eins und Zwei zusammen und brachte stockend hervor: „Ro... Rosanna! Sanna! Mei... meine Sonne!“

Eigentlich änderte sich nichts, denn noch immer war Leander von allen Seite von Shans Schatten umgeben. Und doch änderte sich auf einmal alles, denn die Schatten drängten nicht mehr bedrohlich auf ihn ein, saugten nicht mehr gierig nach seinen Seelenkräften. Nein, die Schatten und die Dunkelheit waren einfach da, ohne aufdringlich zu sein. Sie schienen zu warten.

„Rosanna...“, setzte Leander erneut an, sich von seiner Vision leitend, „Du... du warst einst ein Mensch. Ein glücklicher Mensch. Ein guter Mensch. Du strahltest wie das Licht der Sonne und die Leute lächelten, wenn sie dich auch nur aus der Ferne erblickten.“

Stille. Dann wurde Leander merklich wärmer, als Shan – Rosannas – Schatten sich zurückzog und ihn auf dem Hüttenboden absetzte.

„Was sagst du?“, erklang die tonlose Stimme der Schattenhexe. Leander wusste, dass sie genau verstanden hatte, was er gestammelt hatte. Er war auf Gold gestoßen.

„Ich dachte bislang immer, du wärst nichts weiter als eine Kreatur der Finsternis, geschaffen, um zu zerstören. Aber dem ist nicht so. Du bist ein Geschenk der Sonne, des Lichts. Die Dunkelheit mag ihren Schatten über dich geworfen haben, der dich hungern und zerstören lässt. Dieser Schatten, dunkler als die Nacht selbst, mag sich vor dich geschoben haben, weil das Weltenlicht nicht da war, um ihn zurückzuhalten. Aber in deinem Kern bist du immer noch das reine Licht des Lebens. Aber in deinem Kern trägst du noch immer die junge Frau von damals in dir. In deinem Kern bist du immer noch Sanna aus den Flusslanden.“

Die Schattenhexe blieb stumm. Leander verstummte ebenfalls, und wagte nicht, sich groß zu rühren.

„Du irrst dich“, sprach Shan dann langsam, „Ich bin nicht die, von der du sprichst. Ich bin der Schatten, der schon immer ein Teil von ihr war und eines Tages endlich entfacht werden konnte.“

Varkur war wütend darüber gewesen, an seine menschliche Herkunft erinnert zu werden. Shan hingegen war ruhig, verdächtig ruhig. Und doch schien sie dem Gedanken gegenüber offen, Leander nicht sofort zu verschlingen. Irgendwo in ihrem Innern gab es einen Teil von Shan, der noch als die gute Sanna von damals erkannt werden wollte. Irgendwo in ihrem Innern gab es ein emotionales Zentrum, an das Leander appellieren könnte.

Leander versuchte es mit einer anderen Taktik.

„Nun, eigentlich ist es egal, wer oder was du es einst warst. Wichtig ist, was du sein wirst. Du trägst den Körper Sannas mit dir. Du könntest sie wieder sein. So sage mir, was ist dein Ziel, Schattenhexe? Willst du alle Seelen Andors verschlingen? Tulgors? Des Barbarenlands? Wenn du so weitermachst, kommst du irgendwann ans Ende der Welt und dein Hunger lässt sich nie mehr stillen. Dann wirst du leiden und dich dir selbst stellen müssen. Hättest du dann nicht lieber vorher eingehalten und einen Platz in dieser Welt gefunden? Wenn du hingegen Mäßigung zeigst und hin und wieder Leben bestehen lässt, kann dieses Leben für neues Leben sorgen und du im Gegenzug so lange satt bleiben, wie du existierst. Es ist noch nicht zu spät, einzuhalten.“

„Das ist alles nicht wahr“, erwiderte Shan ruhig, „Du weißt nicht, ob die Welt einen Rand hat. Ich werde mich nie mit meiner Vergangenheit auseinandersetzen müssen, wenn ich nicht will. Und es ist zu spät, um dem Schatten in mir Einhalt zu gebieten.“

Leander blieb stumm. Ihm fiel auf, dass Shan vom Schatten sprach, als wäre er doch nicht ganz Teil von ihr. Und immer noch griff sie ihn nicht wieder an. Was sollte er nun tun?

„Magst du mir noch einmal von Rosanna erzählen?“, erklang nun Shans geisterhafte Stimme, beinahe vorsichtig.

Leander nickte erleichtert und meinte dann: „Was ich in meiner Vision sah, war ein kleines, glückliches Mädchen, welches in den Landen nicht weit von hier lebte. Ich sah das Glück, dass es in die Welt setzte, und ich sah den Stolz in den Augen ihrer Eltern, als sie an Sannas Seite standen. Sanna war besonders, und nicht wegen des Fluchs, den die Dunkelheit auf sie gelegt hatte, sondern aufgrund der Liebe, die sie der gesamten Welt entgegenbrachte. Sanna mag gelitten haben, und der Schatten mag viel Schaden in ihrem Namen angerichtet haben, ja, der Schatten mag sogar ein Teil von ihr geworden sein. Aber noch ist ihre Geschichte nicht zu Ende.“

Eine Zeit lang blieb es wieder still, und Leander vermochte kaum zu hoffen, dass Shan ihn in Ruhe lassen würde. Da sprach die Schattenhexe: „Du sprichst, um dein Leben zu bewahren, o Seher. Das respektiere ich. Und doch sorgen deine Worte für ein warmes Gefühl in meinem Innern. Ein Funken in der Eiseskälte. Es ist ungewohnt, aber nicht unangenehm. Ich verstehe selbst nicht, warum, aber das muss ich auch nicht. Deine Worte spenden Trost, und das tut gut. Eines Tages werde ich deine Seele in mich aufnehmen. Aber noch nicht jetzt. Gehe von hinnen, Seher. Wenn es dich danach dürstet, den Fluch zu lösen, der dich markiert hast, so mögest du dies irgendwann tun. Aber wenn du dich ihm hingeben willst, dann werde ich da sein, um dich mit deiner neugefundenen Macht zu leiten.“

Diese Einstellung erinnerte Leander schon eher an die Shan, die er vor all diesen Jahren in der Rietburg angetroffen hatte. Er atmete tief durch und sprach: „Danke, Sanna. Ich bin auf dem besten Weg dazu, den Fluch zu brechen. Einzig Varkur muss dafür fallen.“

Ein raues Lachen ertönte. Zuerst vermutete Leander, dass es Rosanna ein weiteres Mal mit Freude erfüllte, ihren alten, vergessenen Namen zu hören. Doch dann fuhr die Schattenhexe fort: „Meinst du, mit dem Ende der Blutlinie wäre ein Fluch schon gebrochen? Du bist naiv. Der Fluch hat seit seiner Entstehung Willenskraft und Magie aus seiner Erschaffer und dessen willensstarken Nachkommen gezapft. Er kann auch alleine noch viele Jahre weiter bestehen und von seiner gespeicherten Kraft zehren. Nein, damit der Fluch gebrochen werden kann, muss sich jemand ihm stellen und seine Reserven aufbrauchen lassen. Ihn überwinden. Direkt.“

Erneut hatte Shan gesprochen, als wäre dies das Selbstverständlichste der Welt. Leanders Erleichterung ob des Überlebens dieser Begegnung mit Shan und ob des Empfangens einer weiteren Vision versiegte. An ihre Stelle traten seine üblichen kalkulierenden Gedankengänge.

Wenn Shan die Wahrheit sprach, dann würde es nicht reichen, Varkur zu töten. Zusätzlich müsste jemand nach Narkon reisen und Varatans Fluch persönlich überwinden, ihn alle seine Reserven zehren lassen. Den erschreckten Berichten der Silberzwerge nach besaß der Fluch die Fähigkeit, all jene, die er berührte, alles Wissen über die Außenwelt zu rauben. Manch ein Silberzwerg hatte seine Kumpanen zu weit in die Silberberge wandern sehen, berührt werden und dann ziellos weiterlaufen, immer weiter in die verfluchte Nebelinsel hinein, bis man sie nicht einmal mehr mit einem Fernrohr mehr erkennen konnte.

Wie konnte man einen Fluch überwinden, ohne sich berühren zu lassen? Indem man ihm davonrannte? Indem man vom Fluch sein Gedächtnis löschen und dann mit vorher niedergeschriebenen Notizen wieder auffrischen ließ? Indem man die Grenze des verfluchten Gebiets Narkons durchbrach?

Leander fiel kaum auf, dass Shan davonschlich. Sein Ziel, Callem zu befreien, war schon wieder weiter in die Ferne gerückt. Wie schon so oft zuvor haderte er mit sich selbst. Sollte er Callem erneut aufzugeben versuchen, sich anderen Projekten zuwenden? Bislang war er früher oder später immer wieder auf dieses eine Ziel zurückgefallen, egal, wie oft er sich einzureden versuchte, dass Callem ihm nichts bedeutete, dass Callem nur eine einzige Person in einer so riesigen Welt war, in der Leander einen Sinn finden konnte.

Gemeinsam in einer lebensfeindlichen Welt aufzuwachsen, konnte zwei Brüder derart zusammenschmieden. Menschliche Emotionen. Nicht einmal Leander war gegen sie gefeit.

Kopfschüttelnd sank er in einen unruhigen Schlummer.\bigskip







Am nächsten Morgen – der wohl schon eher gegen Mittag ging – schälte sich Leander nur mit Mühe aus seinem eingekrachten Bett. Sein ganzer Körper schmerzte von der Konfrontation mit gleich zwei übermächtigen Wesenheiten. Varkur und Shan hatten ihn übel zugerichtet. Ihn und seine Hütte. Seine Möbel lagen nicht dort, wo sie sollten, von seinen Pülverchen und Salben ganz zu schweigen. Und sein Geist fühlte sich schwammig und träge an.

Leanders Gedanken drehten sich weiterhin um Varkur und Narkon, ohne dass er dabei auf einen grünen Zweig käme. Er kam bloß zur frustrierten Feststellung, dass diese emotionalen Gedanken an seinen verschollenen Bruder ihn einfach nicht loslassen wollten.

Wie so oft wünschte Leander sich, einen Hinweis aus der Zukunft zu erhalten, welcher ihm einfach mitteilte, wie er Callem befreien könnte. Dieser Wunsch war einst gar einer der Gründe gewesen, warum Leander die Sprache der Zukunft überhaupt erst erlernt hatte.

Zwei Visionen hatte Leander in den letzten Monaten empfangen, bloß zwei. Die erste, nachdem Nomions Geist ihn besessen hatte, als einige eifrige Agren ihn wegen der verletzten Stammesältesten angegriffen hatten. Die zweite gerade erst gestern, als Shan ihm fast das Leben ausgehaucht hatte. Konnte es sein, dass seine stärkeren Visionen durch Kämpfe ausgelöst wurden? Oder durch lebensgefährliche Situationen?

Leander wusste natürlich, dass nicht alle lebensbedrohlichen Lagen in einer Vision resultierten. Varkurs Angriff hatte keine zur Folge gehabt. Und Leander hatte auch schon Vision in alltäglichen Momenten gehabt. Oder?

Jetzt, wo er so darüber nachdachte, war er sich plötzlich nicht mehr sicher. Er konnte sich zumindest nicht daran erinnern, je eine Nachricht aus der Zukunft empfangen zu haben, während er still in seinem Schaukelstuhl gesessen und vor sich hin sinniert hatte. Konnte es sein, dass sein Körper angeregt sein musste, um empfänglich gegenüber der Stimme der Zukunft zu sein? Waren Visionen umso wahrscheinlicher, je heikler seine Situation war, je mehr er sich fürchtete? Konnte das der Grund sein, warum die Visionen abgenommen hatten und unvorhersehbarer geworden waren, je mehr Leander die Sprache der Zukunft erlernt hatte, je natürlicher ihm diese verbotene Kunst erschienen waren?

Eine gewagte These, kein Zweifel.

Doch führte sie in Leanders Kopf zur Bildung eines Plans. Wenn er wahrlich wahrscheinlicher Visionen in Momenten der Gefahr erfuhr, so konnte er sie vielleicht mit Absicht auslösen, indem er sich einer tödlichen Bedrohung stellte. Vielleicht konnte er gar den Inhalt der Vision beeinflussen. Dies würde ihm erlauben, seine Vermutung auf die Probe zu stellen, und im Erfolgsfall zugleich endlich etwas darüber zu erfahren, wie er Varatans Fluch brechen und Callem befreien könnte. Seine letzten beiden Visionen hatten jeweils in Verbindung gestanden zur Gefahr, die ihn in diesem Augenblick bedroht hatte. Sollte Leander vielleicht eine Gefahr in Verbindung zu Varatans Fluch aufsuchen? Zu Callem?

Leander verwarf die Idee, Silberland zu besuchen und die Silberberge zu erklimmen. Dies wäre eine direkte Konfrontation mit dem Fluch, aber eine, die zu hohe Risiken barg. Niemand wusste genau, wie der Fluch agierte. Lieber würde sich Leander einer Gefahr aussetzen, die er genau einschätzen konnte und wusste, wie er sie zu neutralisieren vermochte.

Da aber kam ihm eine andere Idee: Vielleicht müsste die Gefahr gar nicht direkt etwas mit Callem oder Varatans Fluch zu tun haben. Er könnte sich auch einfach in Gefahr bringen und sich vornehmen, sich erst zu retten, sobald er wusste, was er tun könnte, um Callem zu befreien.

Um die Gefahr zu neutralisieren, würde die Stimme der Zukunft ihm dann etwas über Callems Befreiung verraten müssen.

Würden sich seine seherischen Kräfte so leicht austricksen lassen?







\newpage
\section{Vergangenes und Zukünftiges}



Reka hatte sich nicht sehr verständnisvoll gezeigt und Leander keine Phiole von Maros Gift überlassen. Am Ende hatte Leander gar heimlich bei Rekas Schülerin Chada vorstellig werden müssen – ein Fakt, der ihn nicht mit Stolz erfüllte. Chada hatte ihm liebend gerne ein bisschen von Maros Gift gezapft und überreicht (natürlich ohne zu wissen, dass dies gegen Rekas Wünsche ging). Immerhin war bei solchen Aufträgen auf die Helden von Andor Verlass.

Nun stand Leander wieder in seiner Hütte und hielt in seinen Händen eine kleine Phiole mit schwarzen Schlangengift. Vypera-Gift, um genau zu sein. Als er und Callem noch klein gewesen waren, hatten sie gerne den Stummen Wald Narkons erkundet. Nun, der aufmüpfige Callem hatte das gerne getan, Leander war ihm meistens nur hinterhergestolpert.

So war es eines Tages gekommen, dass Leander auf eine Vypera gestoßen war und sich einen Biss eingeholt hatte. Einzig Callems rasches Eingreifen hatte sein Leben damals gerettet – sein Bruder hatte die Schlange mit einem einzigen Streich seines Dolchs aus dieser Welt befördert und sich dann rasch daran gemacht, Leanders Bein abzuschnüren.

„Nicht aufgeben, Leander. Guck mich an. Bleib bei mir!“, ertönte die ruhige Stimme Callems aus seinen Erinnerungen. Leander versuchte, nicht daran zu denken, wie es Callem wohl ging. Was, wenn ihn in der Zwischenzeit eine Vypera erwischt hatte und Callem sich nicht mehr daran erinnerte, wie man mit einer solchen Wunde umging? Hatte Varatans Fluch Callem und dessen Crew schon in den Wahnsinn getrieben? Waren sie überhaupt noch am Leben? Leander verwarf die aufblubbernden Sorgen gekonnt, wie er es schon so oft getan hatte, und wandte sich wieder dem Gift von Rekas silberner Hausschlange zu.

Vypera-Gift wirkte paralysierend und konnte zum Tod führen. „Schwarzes Eis“ wurde es in den zwielichtigen Kreisen des Nordens genannt, und nur ein rechtzeitig eingenommenes Gegengift konnte einen vor größeren Schäden bewahren. Leander hatte sich fest vorgenommen, das Gegengift nicht einzunehmen, ehe er erfuhr, wie er Callem befreien konnte.

Das Gegengift stand ein bisschen abseits in einem Flakon, den Leander einem fahrenden Händler namens Naraven abgekauft hatte. Trotz dessen Anwesenheit war das Unterfangen, das Leander plante, ein äußerst gefährliches. Das schien ja von Nöten zu sein, um seine Visionen zu triggern. Es konnte nur hoffen, dass er, wenn er sich nun Vypera-Gift einverleibte, eine Vision Callems empfangen würde. Es mochte eine exakte Wissenschaft gewesen sein, die Sprache der Zukunft zu erlernen, aber die Visionen der Zukunft auszulösen, das konnte er mit seinen aktuellen Kenntnissen noch nicht in genaue Regeln gießen. Er erinnerte sich daran, wie er sich einst mit Reka darüber gestritten hatte, wie schwammig die Lehre des Zukunftlesens doch war. Vielleicht hatte doch ein Ansatz der Wahrheit in den Worten der Kräuterhexe gelegen.

Sorgfältig zückte Leander seinen Dolch, wägte ein letztes Mal ab zwischen den Risiken seines Vorhabens, seinem Vertrauen in dessen Gelingen und seiner emotionalen Verbindung zu Callem. Dann atmete er tief durch, nannte sich selbst einen Narren und fuhr den Dolch durch seinen linken Handrücken. Diese Extremität könnte er wohl am problemlosesten verlieren, zum Laufen brauchte er beide Beine und rechts war seine dominante Seite.

Die Wunde brannte in der kühlen Abendluft. Leander stellte sich vor, wie sein blaues Blut auf den Boden tropfte. Alsbald steckt Leander sich einen Lederriemen in den Mund, biss darauf und ließ sorgfältig einige Tropfen des schwarzen Vypera-Gifts in die offene Wunde tropfen.

Normalerweise war Leander sehr geschickt darin, Schmerzen zu unterdrücken, aber das Vypera-Gift bot ihm diesbezüglich keine schlechte Herausforderung. Hinzu kam noch, dass Leander beim Kontakt mit dem Gift der Vorfall aus seiner und Callems Kindheit wieder vors innere Auge geführt wurde, was alles andere als angenehm war. Ein unterdrücktes Stöhnen entfuhr seiner Kehle.

Schon breitete sich ein Gefühl der Steifheit in Leanders Hand aus und strahlte in seinen Arm.

Dann wurde alles farbig.\bigskip







Es war kalt. Und es war laut. Schnee und Wind peitschten in Leanders Gesicht, zerrten seine Kapuze von seinem Kopf und ließen sein langes schwarzes Haar umherwehen. Haar, das er schon seit Jahrzehnten nicht mehr besaß.

Er bibberte. Er schien sich in Hadria zu befinden, dem magischen Land von Eis und Schnee. Aus der Ferne war ein animalisches Röhren zu hören. Schreie und Rufe. Meereskreaturen, die dem aufgewühlten hadrischen Meer entstiegen. Nachtgors, die sich an den Schafherden und Einwohnern Hadrias weideten. Ein Trauerfest.

Doch nicht nur das Gebrüll der Kreaturen und die Hilferufe der Hadrier vernahm Leander. Nein, weiter draußen im Meer, so weit draußen, dass man es kaum mehr durch die Schneeschauer erblicken konnte, zuckten Blitze aus dem Himmel und erhellten die Silhouetten mehrerer gut gerüsteter Schlachtschiffe. Von dort waren weitere, eigenartigere Geräusche zu vernehmen:

Der Wind heulte mit einer Stimme, die Flüche vom Himmel schleuderte.

Das Meer flüsterte mit einer zischenden Stimme, die den Kreaturen Schwachstellen an den Schiffen verriet.

Und aus den Tiefen des Abyss erklang ein dumpfes Tröten, welches Leander gehofft hatte, nie wieder in seinem Leben zu vernehmen.

Arkteron, Kenvilar und Oktohan.

Die drei Mächte des Meeres waren entfesselt worden.

Leanders Vision zeigte ihm Hadria zu jener Zeit, als Seekönig Varatan die Mächte herausgefordert hatte. Seine Seefeste Klippenwacht war damals zu einer Ruine degradiert worden. Sein Volk hatte nach Werftheim flüchten müssen, auch wenn diese Nebelinsel damals noch nicht so geheißen hatte. Und die Heere der Mächte des Meeres hatten ihren Blick auf den Ort gerichtet, von dem Varatan seine magischen Waffen erhalten hatte. Hadria wäre in jenen Tagen beinahe überrannt worden. Und auch wenn damals das grausame Opfer Orweyns die Mächte hatte besänftigen können, hatten die Mächte Schnee und Eis mit sich gebracht, welches Hadria bis an den heutigen Tag noch bedeckten.

Warum erinnerte diese Vision ihn an jene schreckliche Zeit? Leander war damals dort gewesen, hatte den Angriff der Mächte mit eigenen Augen gesehen. Leander mochte es nicht, an diese Zeit zurückzudenken. Es war ein Gemetzel gewesen, ein sinnloses Gemetzel.

Orweyn, der mächtigste aller Zauberer, der einst den drei magischen Waffen in der hadrischen Unterwelt ihre Kraft verliehen hatte, hatte als erster erkannt, dass Varatans hochmütige Schlacht gegen die Mächte des Meeres nicht mehr gewonnen werden konnte. Er hatte als erster erkannt, dass die Dunkle Magie, die den magischen Waffen ihre Kraft verlieh, stets ihren Preis forderte. Und so hatte Orweyn sich vorgenommen, die magischen Waffen und jeden, der von der Dunklen Magie wusste, aus dieser Welt zu schaffen. Er hatte in Windeseile den Eisernen Turm erbauen lassen. Und kaum hatte dieser magische Turm über der Feste von Yra gethront, hatte Orweyn eine Versammlung der Festenbewohner einberufen, die das Schicksal Hadrias maßgeblich bestimmen sollte.

Leander erinnerte sich an eine andere Versammlung, die zeitgleich stattgefunden hatte. Er war als interessierter Gelehrter nach Hadria gezogen, nachdem Callem sich der finsteren Kenvilar und der Piraterie verschrieben hatte und Leanders Familie in Varatanien in Verruf geraten war. Leander hatte auch nur als interessierter Gelehrter der Versammlung einiger Magier in Nordgard beigewohnt, die sich ob der nicht abschwächen wollenden Angriffe der Kreaturen zu beraten ersuchten. Mit einem Krachen war Orweyn in der Runde erschienen, seinen grün schimmernden Hammer in der Hand und ein wirres Glitzern in den Augen.

„Die Mächte des Meeres müssen mit einem Opfer besänftigt werden!“, hatte er gegen den heulenden Wind aus dem Schlund Arkterons gebrüllt, „Ein jeder, der von der Dunklen Magie Kenntnisse besitzt, möge mir folgen! Die Magischen Waffen müssen weggesperrt und unser Wissen über sie getilgt werden. Wir versammeln uns bei der Feste von Yra und beenden diese Perversion der Natur ein für alle Mal!“

Eine junge Zauberin war aufgestanden und hatte protestiert. Orweyn hatte nicht lange gefackelt und seinen Hammer auf die Adeptin geschleudert. Leander hatte von Glück reden können, dem ausbrechenden tödlichen Getümmel lebendig entkommen zu sein.

Das nächste Mal, als er Orweyn erblickt hatte, hatte dieser sich eine üble Wunde am Bein zugezogen, aber zusätzlich zu seinem Hammer auch Varatans Helm errungen und eine Handvoll ähnlich denkender (oder aus reiner Furcht mit ihm mitziehender) Zauberer um sich geschart. Gemeinsam war diese Horde durch Nordhom gezogen, hatte Häuser gestürmt und jeden Zauberer gemordet, der sich ihnen nicht anschloss. Sie richteten insgesamt mehr Schaden an als die Armeen von Kreaturen, denn letztere hatten sich zu diesem Zeitpunkt größtenteils ins Meer zurückgezogen, wo die finstere Kenvilar verblüfft zugesehen hatte, wie Orweyn ihre Arbeit für sie tat.

Leander hasste Kenvilar dafür, seinen Bruder auf ihre Seite gelockt und dann im Stich gelassen zu haben. Doch noch mehr hasste er Varatan dafür, den Fluch über Narkon gesprochen zu haben. Wie er Orweyn mit Varatans Helm auf dem Kopf erblickt hatte, hatte er gehofft, dass Varatan inzwischen gefallen sei. Doch dem war nicht so gewesen. Varatan hatte den Untergang seines Flaggschiffs, der Tambur, irgendwie überlebt und sich durch das eisige Wasser des hadrischen Meers an die hadrische Küste geschleppt. Und Kenvilar, die Tückische, hatte von ihm abgelassen.

Damals hatte Leander sich gefragt, warum Kenvilar Varatan nicht nachgetragen hatte. Eine ihrer Töchter war genau wie Callem Teil der Besatzung der Schwarzen Kogge und nun schon seit Generationen auf Narkon gefangen. Hätte Kenvilar Varatan damals den Rest gegeben, so wäre die Besatzung und ihre Tochter vermutlich schon längst wieder auf freiem Fuß.

Heute hingegen fragte sich Leander, ob Kenvilar damals geahnt hatte, welche ungeahnten Ausmaße von Tod und Verderben Varatans Brut in diese Welt setzen würde. Vielleicht war die Menge an Pein, die der Dunkle Magier Varkur in diesen Landen verbreiten würde, in ihren Augen wertvoller als die Freiheit ihrer Tochter? Wer konnte schon wissen, ob ihre Tochter ihr überhaupt etwas bedeutete? Wer konnte schon wissen, ob sie überhaupt wollte, dass die Magischen Waffen unter Verschluss blieben? Vielleicht hätte sie es geliebt, wenn ein mächtiger Zauberer mithilfe der Waffen zu einer vierten Macht des Meeres aufgestiegen wäre.

Orweyn war sich jedenfalls sicher gewesen, dass die einzige Lösung für die drohende Gefahr darin bestand, die Magischen Waffen und alle überlebenden Dunklen Magier im einzig dafür erbauten Eisernen Turm einzuschließen. Nur einer hatte sich vor seinem Mob verstecken können. Zoren hieß er, ein einsamer Einzelgänger, welcher sich bereits seit einem gewaltigen Streit mit dem großen Orweyn von der Zauberergemeinschaft zurückgezogen und in seiner Spitzburg verschanzt hatte. Paranoid, wie er war, hatte er sich beim ersten Anzeichen von Trubel an einem anderen geheimen Ort versteckt, wo selbst die Zauber seiner einstigen Kollegen ihn nicht hatten aufspüren können. Leander hatte mit angesehen, wie Zoren in der Ruhe nach dem Sturm fassungslos durch die Gegend gestolpert war. Anschließend hatte er sich in seine Spitzburg zurückgezogen und war der Legende nach seither nie mehr hervorgetreten. Was er wohl den lieben langen Tag so trieb?

Leander hatte so lange in seinen Erinnerungen geschwelgt, dass ihm erst jetzt wieder bewusst wurde, dass er sich immer noch in einer Vision befand. Einer Vision, die durch Vypera-Gift induziert worden war. Wie lange war er bereits hier? Wie erging es seinem Arm? Er musste weiterziehen.

Die Vision klärte sich ein wenig und Leander Blick fiel auf den Innenhof der Feste von Yra. Niedere Zauberer rannten schreiend umher und trauerten über den Leichnamen ihrer gefallenen Kameraden. Dort drüben umarmte der junge Koraph schluchzend den Körper eines jungen Mannes und versuchte verzweifelt, eine Wunde an seiner Seite zu schließen, die den Zauberer schon längst aus dem Reich der Lebenden befördert hatte. Weit oben in den Himmel ragte der Eiserne Turm, und ein türkises Glühen umgab ihn.

Da oben stand Orweyn, vor dem Eingang zum Turm, seinen Hammer der Stärke locker an seiner Seite haltend, Varatans Helm der Macht auf dem Kopf und Varlion das Flammenschwert hoch in den Himmel gereckt. Er schubste gerade Hombudt, den Zauberer des Raums, ins Innere des Eisernen Turms.

Neben Orweyn rollte ein Braumeister der Feste von Yra ein riesiges Eichenfass ins Innere, aus dessen Seite einige Eiszapfen stachen. Über was für einen sagenumwobenen Inhalt mochte dieses Fass wohl verfügen?

Leander wusste genau, was als nächstes geschehen würde: Orweyn würde gen Himmel blicken und seine letzte Prophezeiung verkünden, mit den Magischen Waffen in den Eisernen Turm steigen und den Turm von innen verschließen, sich und alle überlebenden Dunklen Magier darin festsetzen und damit dem sicheren Hungertod überlassen. Der Turm würde verschlossen bleiben, bis sein Eingang Jahrzehnte später durch Varkur gesprengt werden würde – und auch dann würden die Magischen Waffen dort drinnen verbleiben.

Wie erwartet, reckte Orweyn seinen Hammer und das Flammenschwert gen Himmel und begann, zu seine Prophezeiung zu intonieren:

„Wenn Feuer und Turm miteinander ringen, das Ende aller naht, denn Qurun wird sie bezwingen.“

Qurun war ein hadrisches Wort für das Ende der Welt, die Personifizierung des Bösen, soviel wusste Leander. Doch nun änderte sich die Vision abrupt. Leander verlor kurzzeitig die Orientierung. Als er wieder aufblickte, sah er immer noch den Innenhof der Feste von Yra, doch aus einen anderen Perspektive. Und er nicht mehr Leander. Er war Orweyn. Er blickte durch die Augenlöcher in Varatans Helm der Macht auf eine gelb gefärbte Welt und hielt die beiden anderen magischen Waffen in seiner linken und rechten Hand. Leander spürte Orweyns Entschlossenheit und irre Gewissheit ebenso wie das Brennen der Wunde an seinem rechten Bein, das ihm jemand zugefügt hatte. Ob es Varatan oder ein elender Zauberer gewesen war, wusste er nicht mehr. Es war auch nicht relevant.

Durch Orweyns innere Augen erhaschte Leander einen Blick auf die Vision, die Orweyn damals empfangen hatte, und plötzlich war Qurun nicht mehr nur ein Ausdruck für das Ende der Welt, nein, er hatte ein Bild dazu vor sich: Eine riesige, schattenhafte Gestalt mit spitzen Klauen, langen Zähnen und Augen wie glühender Glut. Gewundene Hörner erhoben sich aus dem schwer im Schatten auszumachenden Schädel, und riesige Flügel trugen das Biest durch die Lüfte.

Unter dem grässlichen Wesen Qurun sah Leander Zauberer des Turms und des Feuers sich gegenseitig bekämpfen, auch wenn Orweyn ihnen diese Begriffe noch nicht zuordnen konnte. Sie schleuderten magische Blitze und Drachenfeuer aufeinander, während über ihren Köpfen Qurun lachte und sich auf dem Eisernen Turm niederließ.

Von Schrecken und Grauen erfüllt, zog sich Orweyn aus dieser düsteren Vision zurück, trennte sich wieder von Leander und verschwand im Innern des Eisernen Turms. Er sprach einige Worte in der Alten Sprache und das riesige Tor fiel krächzend ins Schloss, oder besser gesagt versiegelte sich magisch, denn ein Schloss gab es an diesem Tor gar nicht. Es war nie gedacht gewesen, dass dieser Eingang je wieder geöffnet wurde.

Die Schreie und das Stöhnen der Zauberer im Turm verstummte augenblicklich, das der Zauberer außerhalb blieb hingegen bestehen. Eine türkise Aura umwogte den Eisernen Turm nun, zog Schnee und Eis mit sich und hob diese Schwaden weit in den Himmel. Und wie diese himmelblaue Wogen in den Himmel getragen wurden, schwächte der Sturm, der über Hadria wogte, leicht ab. Die Mächte des Meeres waren besänftigt worden.

Doch Leander achtete nicht darauf. Nein, Leander sah vor seinem inneren Auge weiterhin die Vision, deren Anfang Orweyn gesehen hatte. Er sah, was Orweyn gesehen hätte, wenn er nur ein klein wenig länger mit dem Versiegeln des Turms gewartet hätte. Leander erblickte einige kleine Gestalten, die sich diesem grässlichen Geschöpf namens Qurun näherten und es zurückdrängten. Dort, dieser kleine grüne Bogenschütze, der Varatans Helm der Macht auf dem Kopf trug und gleich drei Pfeile auf einmal im schattigen Leib versenkte, das musste Chada aus dem Wachsamen Wald sein. Dieser Zauberer des Turms, der Varlion das Flammenschwert auf Qurun richtete und einen mächtigen Feuerstrahl in dessen Brust beschwort, das war Eara aus dem Hohen Norden. Oder der Krieger mit dem kalten Blick, der Orweyns Hammer der Stärke auf das Schattenwesen schleuderte, das war höchstwahrscheinlich Thorn aus dem Rietland. Die Helden von Andor, hier in Hadria, mit den Magischen Waffen gegen Qurun kämpfend? Kurz war Leander überrascht, dann lachte er auf. Nun, wer wäre schon so todesmutig, wenn nicht sie? Und wer hatte bessere Chancen, sich Qurun zu stellen?

Dann flackerte die Vision erneut etwas und zeigte Leander das Ende des Kampfes. Qurun sank besiegt zu Boden, die Schatten versiegten und anstelle des riesigen Ungetüms lag nur noch eine kleine Gestalt in einem grauen zerrissenen Gewand im kalten Schnee.

Leander erkannte im grässlichen Geschöpf namens Qurun den ältesten Widersacher der Helden von Andor.

Varkur, der Dunkle Magier.

Er war Qurun gewesen.

Und die Helden hatten ihn besiegt.

Varkur richtete sich mit letzter Anstrengung aus, spuckte Blut und rief etwas in den heulenden Wind. Dann sank Varkur, Sohn des Varkmar, in dunklen Nebel und schied aus diesem Leben.

Leanders sehendes Auge verschloss sich.\bigskip







Nur langsam kam Leander wieder zu Besinnung. Als erstes merkte er, dass er auf seinem Hüttenboden lag und ihm der bittere Geruch nach Sternkraut in sein Riechorgan stach. Dann erinnerte er sich an seinen Plan, eine Vision über Callem zu erhalten. Der Plan hatte funktioniert. Das schwarze Gift hatte es tatsächlich geschafft, eine Vision zu triggern! Eine ungewöhnlich lebhafte und ungewöhnlich lange war es gewesen, und erneut hatte er in die Vergangenheit statt in die Zukunft geblickt.

Erschrocken wurde ihm bewusst, dass er zwar Maros Gift, nicht aber das Gegenmittel eingenommen hatte. Wie viel Zeit war vergangen? War seine Hand noch zu retten? Rasch versuchte er, sich aufzurichten, da traf ihn eine Hand unsanft auf der Brust und hinderte ihn daran.

„Ts, Ts, Ts“, erklang ein kritisches Schnalzen, „Wenn ich dir deinen Wunsch nach Schlangengift verweigere, hat das einen Grund. Das heißt nicht, dass du stattdessen einfach mein Mädel danach fragen sollst.“

Reka. Natürlich hatte Reka ihn gefunden.

Leander sank seufzend auf den Boden zurück. Erleichtert, dass sein Leben nicht mehr in Gefahr war. Aber auch furchtsam vor den Worten der Hexe. Und diese war noch nicht fertig mit ihm. Sie sprach weiter, und zum ersten Mal konnte Leander nicht einmal den Ansatz von Humor in ihrer Stimme erkennen:

„Du kannst von Glück sprechen, dass Chada mir rechtzeitig davon berichtet hat. Wenn ich einige Minuten später hier aufgekreuzt wäre, hätte ich dich nicht wieder zurückbringen können.“

„Ich habe doch gar nichts gesagt“, erklang die protestierende Stimme Chadas aus dem Hintergrund. Sie erinnerte Leander an die Vision, die er soeben erblickt hatte. Chada und die restlichen Helden würden Qurun besiegen. Und Qurun würde Varkur sein. Interessante Informationen, zweifelsohne, insbesondere da sie ihm erlaubten, sich keinen ausgeklügelten Mordplan für Varkur ausdenken zu müssen. Egal, was er tun würde, Varkur würde nicht ewig am Leben bleiben. Dennoch hatte ihm die Vision nicht das gezeigt, was er sich gewünscht hatte. Nichts zu Callem und wie man ihn von Varatans Fluch befreien könnte. Ganz hatte er seine seherischen Kräfte also nicht austricksen können.

„Hast du etwas zu deiner Verteidigung zu sagen?“, sprach Reka. Leander versuchte, nicht daran zu denken, dass er entgegen ihres Wunsches Meres nach Narkon geschickt hatte, und entgegnete stattdessen:

„Wenn ich dich so sehr störe, warum hast du mich überhaupt gerettet?“

Reka blieb einen Moment stumm. Dann meinte sie: „Nenn mich verrückt, Leander, aber ich habe das Gefühl, dass du deine große Rolle in den Legenden von Andor noch nicht gespielt hast. Ich weiß, dass du viel Übles angerichtet hast in deiner langen Lebzeit. Aber ich glaube, dass du auch noch viel Gutes anrichten wirst.“

Leander gluckste schwach auf. Dafür konnte er die Hexe wirklich verrückt nennen.

Diesmal hatte sie auch keinen Eintopf für ihn übrig, als sie ihn verließ. Nein, Leander musste ihr danken dafür, dass sie sein Leben gerettet hatte. Und das tat er auch. Anschließend dirigierte er sie und Chada aber auch schon wieder unwirsch aus seiner kleinen Hütte hinaus und sortierte seine Gedanken im Angesicht seiner jüngsten Vision.

Er wusste nun, wie Varkur sterben würde. In Hadria. Durch die Hand der Helden, wie könnte es anders sein? Nur wie bald? Die Helden wirkten nicht so, als würden sie Andor so bald verlassen wollen. Varkur könnte sie natürlich in den Norden locken. Aber warum würde er das tun?

Leander studierte und knobelte an diesem Rätsel herum, bis ihm endlich ein Ansatz einfiel. Die Krahder aus dem Süden. Das Riesenvolk hatte vermutlich keine Ahnung davon, dass Tarok gefallen war, dass der Weg nach Andor für sie offen stand, ja, hatte vermutlich nicht einmal eine Ahnung davon, dass da im Norden ein Königreich lag, welches von aus Krahd geflohenen Ambacus gegründet wurde, und dass Mitglieder dieses Königsreichs einen der ihren erschlagen hatten. Sobald sie davon erfuhren, würden sie wohl mit Vergnügen Rache nehmen wollen.

Die meisten Krahder waren nun mal Feiglinge. Folglich müsste jemand die furchterregenden Helden von Andor aus dem Weg schaffen, damit sie in Andor einfielen. Und jemand müsste die Krahder über den freien Weg in den Norden informieren. Diese Rollen könnte Varkur übernehmen. Und dies gäbe Varkur einen guten Grund, in den Norden zu reisen und die Helden hinter sich herzulocken.

Kurz verzog Leander seine Mine, als er an die vielen Bauern der Andori denken musste, von denen bestimmt einige bei einem solchen Überfall der Krahder umkommen würden. Er unterdrückte die Gedanken. Die Bauern bedeuteten ihm nichts, und Krahder in Andor wären ganz nebenbei keine schlechte Gelegenheit für Leander, dunkles Wissen zu ergattern. Aber der Angriff der Krahder wäre nicht einmal das wichtigste, sondern nur ein Grund für Varkur, in den Norden zu ziehen. Wichtiger war, dass sie Helden von Andor Varkur in den Norden folgten und dem Dunklen Magier dort ein Ende bereiteten. Dass sie es möglich machten, dass Varatans Fluch gebrochen werden konnte.

Und wenn die Helden erst einmal im Norden fertig waren, würden sie zurückkehren wollen, vermutlich in Eile, wenn es in Andor dann nicht allzu rosig zu und her ginge. Und wenn ihr Schiff auf Narkon kenterte... Könnte Leander einen Pakt mit einer Macht des Meeres eingehen? Könnte er Arkteron, Kenvilar oder Oktohan dazu bewegen, das Schiff der Helden in die Stachelklippen Narkons zu dirigieren? Wer könnte schon mit Sicherheit Varatans Fluch die letzten Kraftreserven aufzehren lassen, wenn nicht die Helden von Andor?

Aus den wirren Gedanken in Leanders Kopf kristallisierte sich immer deutlicher ein grandioser Plan heraus. Erschöpft sank er in seinen Schaukelstuhl und atmete schwer durch. Es würde riskant werden. Aber um Callem zu befreien, hatte er schon viel Riskanteres durchgezogen.

Leander massierte seine Schläfen. Shan war vermutlich gerade auf der Suche nach Varkur, um dabei zu sein, wenn dieser Hademars Seele aus dem Drachenherzen befreite. Und Shan konnte die magischen Spuren ihres Schattens erheblich schlechter verbergen als Varkur die seinen. Wenn er Shans Spur folgte, würde er höchstwahrscheinlich auf Varkur stoßen können.

Danach müsste Leander seinen Plan nur noch Varkur unterbreiten. Nun, zumindest den Teil, der nicht Varkurs Tod beinhaltete. Am besten so, dass der Dunkle Magier das Gefühl hatte, von selbst drauf gekommen zu sein.

Entschlossen griff Leander nach seinem Stab.\bigskip







Varkur, der Dunkle Magier, halb Mensch, halb dunkler Nebel, erreichte das Graue Gebirge. Langsam schälte sich sein Körper aus dem Nebel. Mit geschlossenen Augen stand er da, der eisige Wind schlug ihm ins Gesicht und der Schnee hüllte ihn mehr und mehr ein. Er war in Hochstimmung.

Einige Wochen hatte er nun schon vergeblich damit verbracht, Hademars Geist aus diesem verfluchten kleinen Steinchen zu befreien versuchen. Sein Kopf schmerzte von all den niederen Kreaturen, in die er Hademars Geist schon projiziert hatte, nur um kurze, unverständliche Antworten zu erhalten.

Hademars Seele, sein Geist, oder wie auch immer man das nennen wollte, hatte sich stets kontaktfreudig gezeigt. Selbst wenn er nicht gerade einen Gorlot steuerte, konnte der Nekromant den schwarzen Schemen unter der roten Edelsteinoberfläche des Drachenherzens mit einigem Geschick verformen. Ohnehin verkraftete Hademar die Zeit in Gefangenschaft überraschend gut. Aber vermutlich war er sich schon daran gewöhnt, körperlos zu sein. All die Jahre, die er seit dem Verlust seines Körpers damit verbracht haben musste, die Überreste seines Geistes langsam zusammenzusetzen und erstarken zu lassen... Varkur wollte sich nicht vorstellen, wie sich das hatte anfühlen müssen.

Varkur vermisste es, Hademar von Gesicht zu Gesicht gegenüberzustehen. Die beiden hatten gemeinsam gute Zeiten erlebt in Hadria. Ehe sie auseinander getrieben worden waren vom Schicksal. Oder besser gesagt von ihren persönlichen Rachefeldzügen und bösartigen Plänen. Dort waren sie nie auf einer Wellenlänge gewesen. Varkur hatte rasch handeln wollen, Kreaturen zu offenen Angriffen nutzen, den Drachen Tarok zähmen, die Völker fremder Länder gegeneinander locken, Könige umbringen und am Ende mit einer riesigen Armee unter seiner Kontrolle gen Hadria ziehen.

Hademar hingegen hatte sich so stark wie nur irgendwie möglich vom für ihn verhassten Andor fernhalten und im Geheimen seine Macht mehren wollen. Alte Schriften studieren, leere Festen erkunden, mächtige Formen der Magie meistern. Erst wenn er sich (fälschlicherweise) sicher gewesen war, dass niemand gegen ihn ankommen könnte, erst dann hatte er sich offenbart.

Am Ende war aus dem Dunklen Duo eine schwache Bekanntschaft geworden. Hin und wieder hatten sie sich Briefe gesendet und den jeweils anderen darüber aufgeklärt, wie viel besser ihre Angehensweise sich für sie bezahlt machte.

Am Ende hatte es für sie beide nicht funktioniert. Die Helden von Andor waren ihnen beiden in die Quere gekommen, immer und immer wieder. Varkurs Gesicht verzerrte sich beim Gedanken an sie zu einer grimmigen Fratze.

Einst waren Hademar und er zwei begeisterte Studenten der geheimen Magieformen in Hadria gewesen. Nun war Varkur ein halber Echsenmensch, umgeben von einem ihn verzehrenden brennenden Nebel, und Hademar war nichts weiter als ein schwarzer Schatten unter der Oberfläche eines kristallenen Runensteins, der hier und dort durch eine besessene Kreatur einige Worte sprechen konnte.

Aber nicht mehr lange. Hademar war es gelungen, Varkur in den letzten Wochen einige konkrete Hinweise zu übermitteln.

„Licht.“ „Mond.“ „Altar.“

Solche Wörter hatten die Gorlots, Wargors und Skrale gekrächzt, ehe die Verbindung gebrochen worden und sie wieder zusammengeklappt waren.

Varkur hatte einen genügend geübten Blick für die fließenden Bänder der Magie, um zu spüren, dass mächtige Drachenmagie Hademar in diesem Stein einschloss. Nach und nach hatte er es mit Hademars Hinweisen geschafft, den Prozess herauszufinden, der dies bewerkstelligt hatte. Das Licht des roten Mondes hatte den Stein in einem Ritual blutrot gefärbt. Ein Alter aus purem Roteisen hatte geholfen, dieses Licht zu kanalisieren. Das bedeutete, dass Zwerge involviert waren. Diese unterirdischen Ratten! Varkur hätte sie schon längst aus ihren engen Gängen ausräuchern sollen.

Heute hatte Hademar Varkur endlich in die unterirdische Kammer gelenkt, wo sein Geist eingesperrt worden war. Tief unter das Graue Gebirge war Varkur geflogen, hinein in die Korn-Schlucht, ins unterirdische Reich der Schildzwerge aus der Zeit, wo sie noch nicht einmal Schildzwerge genannt worden waren, nach links, dann zweimal rechts, nach oben und dann noch weiter nach unten, stets dem schwarzen Pfeil folgend, den Hademars schwarze Seele im Edelstein für ihn formte.

Nun stand er in einem kreisrunden Raum. Dieser war leer bis auf den Blutstein-Altar in seiner Mitte und eine riesige Gestalt am gegenüberliegenden Ende des Raumes. Dort lag ein versteinerter Drache, die Überreste des einst so edlen Kardòl, den rissigen Kopf in seinen stacheligen Schwanz verbissen. Aber Varkur war nicht seinetwegen gekommen.

Die Spuren des magischen Rituals waren nicht mehr frisch, aber dennoch deutlich zu erkennen. Hier war Hademar im Drachenherz eingeschlossen werden. Schlampige Arbeit, aber dennoch schwer zu durchbrechen. Varkur schüttelte seinen Kopf und platzierte das Drachenherz vorsichtig auf dem großen Alter aus purem Roteisen. Blutsteinaltar hatten ihn die Zwerge genannt, nach dem Blut der gefallenen Drachen, welches dem Stein seine Eigenschaften verliehen hatte.

Der rote Mond schien nicht, durch den gezackten Spalt in der Decke konnte Varkur nur den klaren Sternenhimmel erkennen. Das Sternbild des Hornfalken, wenn er sich nicht irrte. Ein gutes Omen. Er brauchte das Licht des roten Mondes nicht, um Hademar zu befreien.

Varkur schlug seine Ärmel zurück und enthüllte echsenhafte Klauen, die er knackend spreizte. Zeit, seinen alten Freund zurückzuholen.

Er ritzte sich seinen Handrücken auf und spritzte sein schwarzes Blut über den Altar. Sein Mund öffnete sich und enthüllte spitze Zähne, als er grinste und ansetzte, die entsprechenden magischen Worte der Alten Sprache auszusprechen, die Hademar befreien würden.

Urplötzlich wurde es kalt um Varkur. Der Dunkle Magier spürte, wie \textit{Etwas} in seinen Geist eindrang und seine Gedanken in Eiseskälte tauchte. Alarmiert stolperte er nach vorne und stieß gegen den Blutsteinaltar.

Sein Körper fing sich von selbst. Von Horror erfüllt, spürte Varkur seine lange Zunge in seiner Mundhöhle umherwandern und seine scharfen Zähne entlanggleiten. Seine eigenen Lippen öffneten sich und er hörte seinen eigenen Mund krächzend sprechen: „Tut mir leid, Hadrier. Ich kann das nicht zulassen.“

Varkur versuchte zurückzuweichen, zum Schlag auszuholen, seinen Stab zu greifen, den er achtlos beiseite geworfen hatte.

Nichts konnte er tun. Alles erfolglos, sein Körper gehorchte ihm nicht mehr.

Einzig der Nebel der Dunklen Magie gehorchte ihm noch. In Varkurs Auftrag wallte dieser durch den Raum, spaltete Stein und brach Felsen, suchte nach einem Gegner, dem er sich stellen konnte. Doch es gab keinen solchen Gegner. Denn der Gegner befand sich in Varkurs Kopf.

„Du verstehst das bestimmt“, sprach sein Mund nun, und wie ein Fremder in seinem eigenen Körper hörte Varkur die pure Finsternis aus sich sprechen, „Hademar ist ein mächtiger Nekromant. Wird er nun befreit, könnte er zur mächtigsten Person dieser Welt werden. Das darf ich nicht zulassen.“

War das die Dunkle Magie, die da aus ihm sprach? Hatte sie endgültig seine Menschlichkeit in sich aufgesogen, ihn endlich endgültig in ihre Gewalt gekriegt? Varkur ließ sich ergeben sinken. Nun, nicht wirklich sinken, er hatte ja keine Kontrolle über seinen Körper mehr. Aber er hörte auf, gegen die fremde Macht anzukämpfen. Die Dunkle Magie brannte schon seit Jahrzehnten in ihm, und je länger er von ihr Gebrauch gemacht hatte, desto öfter hatte sich dieses zweischneidige Schwert gegen ihn gewandt und ihm mehr seiner Menschlichkeit genommen. Er war dazu verdammt, früher oder später als ausgebrannte Hülle der Dunklen Magie zu dienen. Dieses Schicksal hätte ihn ohnehin irgendwann erreicht. Nur schade, dass es so viel früher als später der Fall gewesen war.

Ergeben nahm Varkur wahr, wie seine Hände sich öffneten und seinem Zauberstab geboten, zu ihm zu springen. Sein Stab tat wie geboten, und die Seelen der unzähligen darin eingefangenen gequälten Geister stöhnten auf, als die finstere Macht in ihre Energie griff und ihre Kraft entfesselte. Dann hoben sich Varkurs Hände, den Zauberstab fest umschlossen. Seine Kehle brach in einem dämonischen Lachanfall aus, als der Stab auf das Drachenherz niederfuhr, fest entschlossen, Hademar ein für alle Mal auszulöschen.

Stattdessen blieb der Stab kaum eine Handbreit über dem Drachenherz in der Luft stehen, als wäre er auf eine unsichtbare Wand getroffen. Varkurs echsenartige Augen rollten ohne seine Führung in ihren Höhlen umher und versuchten, im Dunkel die Ursache dieser Störung zu erkennen. Als wäre dies überhaupt möglich gewesen, so wurde der Raum noch kälter, und vor allem eines, dunkler. Kein Stern war mehr durch die Ritze in der Decke zu erkennen.

Immer finsterere Schatten flossen in den Raum und Varkur erkannte, dass das, was seinen Stab davon abgehalten hatte, das Drachenherz zu zersplittern, ein finsterer Tentakel aus purem Schatten war, der aus einem schattigen Gebilde zu strömen schien, welches sich soeben in den Raum zwängte.

Zwei große gelbe Augen leuchteten in dem Gebilde auf und die Schatten formten sich zu einem dunklen Schemen, der entfernt an die Gestalt einer Frau erinnerte. Varkur wäre zurückgezuckt, hätte er können, so sehr erinnerte ihn diese Erscheinung an seine N...

„Es ist sie nicht“, sprach er in Gedanken zu sich selbst, „Es kann nicht sie sein.“

Da sandte der Dunkle Magier seinen Nebel aus, der schwarz und bissig wie ein Tier auf die Schattenfrau zuglitt und sie umwallte. Oder vielmehr versuchte, sie zu umwallen. Die Schattententakel der Schattenhexe wogten und verformten sich zu einem breiten Schild aus purer Dunkelheit, welcher sich zwischen die Frau und seinen Nebel stellte und diesen zurückdrängte. Wie war das möglich?

Varkurs Körper schnellte nach vorne, als er von der finsteren Macht, die ihn kontrollierte, zum direkten Angriff dirigiert wurde. Varkurs Krallenhände stürzten sich auf den Schattenschild und durchbrachen dessen Oberfläche, sorgten für einen Spalt, durch den sein Nebel zur Hexe dringen konnte. Sie schrie nicht auf, aber ihre ganze Form zitterte und verschwamm. Gemeinsam stürzten sich Varkur und die Macht, die seinen Körper führte, auf die Schattenhexe.

Nun, da er sie genauer erkennen konnte, atmete Varkur erleichtert auf. Das war nicht N... Ni... es war sie nicht.

Die Schatten der Schattenhexe zerfaserten und flossen um Varkur herum, dessen Körper herumschleuderte und die Schattenhexe zu fassen versuchte.

„Narr“, drang eine flüsternde und doch durchdringende Stimme aus dem Schatten, „Ich bin auf deiner Seite, Dunkler Magier. Ich will genau wie du den Nekromanten aus dem Stein befreien.“

Schwarze Schemen flossen an Varkurs Körper hoch und verfestigten sich um seine Hand- und Fußgelenke. Mit übermenschlicher Kraft zerrte die Schattenhexe an ihm und bugsierte ihn wieder in die Nähe des Drachenherzens, Hademars Edelstein, der noch immer auf dem Blutsteinaltar lag.

Varkur ließ seinen Dunklen Nebel abziehen. Wenn diese Schattenhexe wirklich Hademar befreien wollte, waren sie wahrlich auf derselben Seite. Er spürte, wie seine Muskeln von der fremden Macht in seinem Körper versteift wurden und sich erfolglos gegen die Kontrolle der Schattenhexe zu wehren versuchte.

Schon hatten Varkurs Klauenhände den Edelstein wieder ergriffen und in die Höhe gehoben. Unruhig nahm Varkur war, wie die Schatten der Schattenhexe schwächer wurden, je näher sie dem Edelstein kamen. Das Licht der Gemme schwächte sie. Nichtsdestotrotz vermochte sie immer noch, seinen Körper zu führen. Ein unterdrückter Aufschrei entfuhr ihr, dann hielt sie den Stein mit Varkurs Händen in die Höhe und begann, einen Satz in der Alten Sprache zu sprechen.

Varkurs Kehle entfuhr ein Aufschrei und Varkurs Magen stellte sich auf den Kopf. Sein Körper sackte zusammen und ein keuchender Husten schüttelte ihn, als eine plumpe, blau-graue Masse aus seinem Körper floss und sich vor ihm zusammensetzte. Dann war es also doch nicht die Dunkle Magie gewesen, die ihn kontrolliert hatte. Nur ein fremder Geist.

Auf dem kalten Steinboden liegend, erkannte Varkur, wie der Geist kurz eine konkrete Gestalt annahm: Eine schwebende Masse von grauen Tüchern, auf der ein kleiner gemeiner Krahder-Kopf saß. Dann richtete der Geist eine siebenfingrige Hand auf die Schattenhexe und brüllte: „Wer wagt es, sich mir zu widersetzen?! Ich bin Nomion, der Meister des Urtrolls! Eine kleine Schattenhexe kann mir nicht in die Quere kommen!“

Nomions Form zerfaserte und stürzte auf die Schattenhexe zu. Er drang mühelos durch ihren Schattenschild und verschwand in der Dunkelheit, die die Schattenfrau umgab. Selbige kreischte auf. Ihre gesamte Gestalt waberte und zitterte, als sie sich gegen die Kontrolle des Krahder-Hexers zu wehren versuchte.

Gut, dachte Varkur, soll er nur versuchen, sie zu kontrollieren. Es kostete Nomion bestimmt einiges an Energie, einen neuen Körper zu übernehmen. Irgendwann müssten seine Reserven aufgebraucht sein. Und die Schattenhexe schien ihm einen angemessenen Kampf zu bieten.

Die ganze Welt drehte sich unschön, als Varkur seinen Körper in die Höhe stemmte. Erleichtert streckte und dehnte er sich. Sein Körper gehörte wieder ihm. Zeit, das zu tun, warum er überhaupt hierhergekommen war.

Varkur ergriff das Drachenherz, hob es in die Höhe und setzte erneut zum Satz in der Alten Sprache an, den Hademar aus seinem Gefängnis befreien würde.

Da wurde ihm von einem Schattententakel das Drachenherz aus der Hand geschleudert. Frustriert wollte er sich umdrehen, da umschlangen bereits weitere Tentakel seine Arme und Beine und zogen ihn nach vorne. Das Gesicht der Schattenhexe erschien vor ihm, doch nun waren ihre Züge noch unmenschlicher und ihre Ohren spitz wie die eines Krahders, als eine viel tiefere, doch immer noch geisterhafte Stimme aus ihrer Kehle erklang:

„Wie schon gesagt: Ich kann das nicht zulassen, Hadrier. Hademar darf nicht befreit werden.“

Offenbar war es Nomion gelungen, den Geist der Schattenhexe zu vergiften, ohne dass sie im Gegenzug den seinen hätte auslöschen können. Nomion schlang einen weiteren Schattententakel um Varkurs Kehle und drückte zu. Varkur hatte selbst gesehen, wie rasch die Absenz von Luft einem Köper das Leben ausquetschen konnte, und haderte nicht lange mit einer Antwort.

Verzweifelt saugte er die Dunkle Magie in sich ein, viel mehr, als er für eine lange Zeit aus ihr geschöpft hatte. Er spürte, wie in dieser lebensbedrohliche Situation ein letzter Damm in ihm brach, wie die Magie durch sein Adern jagte. Und wie sie ihm erlaubte, die Schatten zu packen, die ihn umgaben.

Mit letzter Kraft griff Varkur nach den Schatten der Schattenhexe, aber nicht mit seinen Klauen, sondern mit der Dunklen Magie. Er zog sie zu sich, leitete sie durch den Blutsteinaltar, verband sie mit der allgegenwärtigen Drachenmagie und versuchte verzweifelt, so viel wie möglich in sich aufzunehmen, um Nomions Angriff widerstehen zu können. Sein schwarzer Nebel drang in die Schatten ein, doch anders als vorher wurde er von diesen nicht abgewehrt. Stattdessen verbanden sich sein Nebel und die Schatten der Schattenhexe. Varkur zog und zerrte an der neuen Verbindung. Für einen kurzen Augenblick löste sich der Druck auf seine Kehle, flossen die Schatten von dem Körper der Schattenhexe weg und auf ihn zu. Varkur glaubte, inmitten der Dunkelheit die Gestalt einer jungen Frau mit einem traurigen Lächeln auf den Lippen zu sehen. Dann ballten sich die Schatten wieder um sie und der Körper der Schattenhexe richtete sich wieder auf, noch größer als zuvor. Der Druck auf Varkurs Kehle verstärkte sich wieder.

Erneut zerrte Varkur an dieser seltsamen Verbindung seiner Dunklen Magie und des Schattens der Schattenhexe, noch stärker als vorher. Diesmal fühlte er, wie etwas vor ihm auseinanderriss, das Ziehen einer anderen Macht am Schatten versiegte und plötzlich unmöglich zu bändigende Mengen an Dunkler und Schattenmagie durch seinen Körper strömten, ihn transformierten.

Varkurs Muskeln schwollen an und seine Haut verschuppte vollends. Ein Prozess, den die Dunkle Magie schon seit Jahrzehnten an ihm vornahm, wurde in Sekundenschnelle ins Extrem getrieben. Vor Varkurs glühenden Augen wuchsen seine Klauenhände auf die zweifache, dreifache, vierfache Größe an. Seine Halsmuskeln spannten sich und schoben den kläglichen Rest von Nomions Schattententakeln mühelos zur Seite. Seine Kopfhaut spannte sich, als seine Hörner in die Höhe schossen, seine Kapuze zerrissen und die Schatten des Raumes durchbrachen.

Varkur fühlte, dass er plötzlich so viel mehr war als sein stetig wachsender Körper und der allgegenwärtige Nebel der Dunklen Magie. Er fühlte den verfluchten Schatten im Raum, wie er über den kalten Steinboden strich und die Dunkelheit liebkoste. Fast war es Varkur, als flüsterte dieser Schatten zu ihm: „Ich bin dein. Nutze mich!“

Probehalber befahl er dem Schatten im Geiste, sich um ihn zu versammeln. Prompt floss der Großteil der magischen Dunkelheit, die diesen Raum unter dem Grauen Gebirge erfüllt hatten, zusammen und umgab den sich transformierenden Körper Varkurs, verband sich mit seiner Schuppenhaut, stärkte seine Klauen, maskierte sein Gesicht.

Der Körper der Schattenhexe waberte weiterhin vor ihm, aber im Vergleich zu den Schatten, die Varkur nun umgaben, sahen die wenigen Fetzen von Dunkelheit, die noch um die Schattenhexe wogten, beinahe lächerlich aus. Varkur hatte der Schattenhexe soeben den größten Teil ihrer Macht entrissen. Die Schattenhexe, oder vielmehr ihr Körper, heulte auf und sank zusammen. Nomions Geist löste sich von ihr und nahm ein weiteres Mal Gestalt an, nur um kreischend das Weite zu suchen.

Die Schattenhexe richtete sich auf und blickte Varkur aus leeren Augen an.

„Wie... was?“, stammelte sie schwach.

Varkurs Körper saugte weiterhin ungehindert Dunkle Magie in sich auf, gestärkt durch die Kraft der Schattenseelen. Seine Klauenhände wurden zu schwer, um aufrecht stehen zu bleiben, und so stürzte Varkur auf sie. Seine unförmigen langen Hörner stießen bereits gegen die Decke des Zwergenraums. Die Schatten, die einst der Schattenhexe gehört hatten, wurden zu riesigen Schwingen, die sich hinter ihm erhoben und den ganzen Raum ausfüllten, ja, gar aus dem Raum hinausquollen. Und als er seine Stimme erhob, um die Schmerzen in seinem Innern auszudrücken, erklang an ihrer Stelle ein tiefes animalisches Röhren.

Varkur fiel kaum auf, dass nun auch die Schattenhexe hastig das Weite suchte. Er drehte und wendete sich, versuchte, im staubigen Steinraum Hademars Edelstein zu finden. Es gelang ihm nicht. Das Drachenherz war verschwunden. Und selbst wenn er es finden könnte, könnte er es in seinem jetzigen Zustand wohl kaum transportieren.

Varkur tauchte in den Strom der Dunklen Magie, der ungehindert in ihn hineinströmte, und versuchte, ihn zu versiegen zu lassen. Er spürte, wie seine Wandlung ein klein wenig stagnierte, aber es war, als würde er versuchen, den Fluss des Tatru mit einem einzelnen Holzstamm aufzuhalten. Erneut stieß Varkur mit seinen Hörnern gegen die Höhlendecke. Wenn diese Transformation so weiterginge, würde er bald hier unter dem Gebirge feststecken. Er musste hier raus, und zwar rasch!

Varkur brach durch die Tür des Raumes in den dahinter liegenden Gang. Ein Paar Gorlots heulten auf und sprangen davon, so schnell sie konnten. Varkur war schneller. Er vermochte sich nicht mehr an den genauen Weg erinnern, den er genommen hatte, um hierher zu kommen. So schrammte er blindlings durch die unterirdischen Gänge des Zwergenreichs, schlug Türen auf und Säulen um, kratzte tiefe Krallenspuren in uralten Stein und schlug mit seinen Hörnern Löcher durch die Höhlendecke. Seine Schattenflügel wogten um ihn herum, wie es einst sein finsterer Nebel getan hatte, zertrümmerten Stein, pulverisierten Fels und brachen Gorlot-Knochen, wo immer sie konnten.

Ein dumpfes Dröhnen erklang, als hinter ihm ein Gang nach dem anderen einstürzte. Varkur legte noch weiter an Tempo zu und holperte und polterte auf allen vieren durch das verlassene Zwergenreich, in der Hoffnung, die freie Luft zu erreichen, ehe sein hünenhafter Echsenkörper zu groß war, um sich fortbewegen zu können.

Er erreichte eine Sackgasse, prallte mit voller Wucht in eine Wand – er spürte es kaum, so sehr federten seine neuen Schatten den Aufprall ab – und hinter ihm stürzten Tonnen von Geröll hinunter, als der Gang zusammenstürzte und den Rückweg blockierte.

Es war stockfinster, nass und kalt. Varkurs Körper schwoll weiterhin an und presste gegen die Höhlenwände. Es wurde immer schwerer, sich noch zu bewegen. Erst als er quasi den ganzen Raum ausfüllte und physisch nicht mehr weiterwachsen konnte, spürte Varkur, wie der Strom der Dunklen Magie in ihn versiegte. Er befand sich nun in einem Raum irgendwo unter dem Grauen Gebirge, der seit Jahrzehnten, vielleicht gar Jahrhunderten von keiner Zwergenseele betreten worden war, eingequetscht in einem unmenschlichen Echsenkörper, umgeben von Dunkler Magie und fremden Schatten, die verzweifelt gegen den von allen Seiten auf ihn drückenden Felsen ankämpften.

Es grauenvolles Knirschen zog sich durch den Fels und Varkur hätte sich seine nun spitzen Riesenohren zugehalten, hätte er sich noch bewegen können. Er fürchtete, das ganze Graue Gebirge würde gleich auf ihn einprasseln und ihn endgültig aus dieser Welt holen. Varkur, der Dunkle Magier, erdrückt von einem Berg. Was für ein lächerliches Ende!

Krachend gab die eine Seitenwand des engen Raums nach, in dem Varkur eingeschlossen war. Varkurs purzelte aus der neu entstandenen Öffnung in Felswand ins Freie (und in den freien Fall), während der kleine Raum hinter ihm kollabierte. Die Höhle hatte direkt an der Korn-Schlucht gelegen und Varkur stürzte unkontrolliert in deren Tiefe!

Varkur ächzte, während der Wind an seinen Schwingen riss, doch konnte er seinen Fall mit den Schattenflügeln aufhalten. Tief unter sich glaubte er, einen riesige steinerne Gestalt mit einem abgebrochenen Horn in der Schlucht liegend zu erkennen. Dann zog sein Blick nach oben und blieb auf der Ruine der Zwergenbrücke Tiefenfall liegen, die die tief darunter ruhende Gestalt der Legende nach vor ihrer Fertigstellung bereits wieder eingerissen hatte. Mit seinen mächtigen Schwingen schleuderte er sich in die Höhe und katapultierte sich aus der Korn-Schlucht heraus.

Erleichtert stellte Varkur fest, dass der Strom der Dunklen Magie versiegt war und seine deformierte Schattengestalt nicht mehr immer größer wurde. Tatsächlich begann sein Körper schon damit, einen Großteil der Dunklen und Schattenmagie freizugeben. Varkur schrumpfte, die Schatten um ihn herum zerflossen, und mehr torkelnd als elegant landete der Dunkle Magier am Rande Tiefenfalls.

„Keine Sorge, wir sind nicht für immer verschwunden“, schien diese ungebändigte Kraft mit den flüsternden Stimmen hunderter Seelen zu wispern, „Wir werden wieder an deiner Seite stehen, wenn du uns wieder brauchst.“

Die fremden Seelen verflogen völlig und hinterließen Varkurs kleinen menschlichen Körper auf dem nackten Fels liegen. Es war eiskalt. Er drehte sich elend zur Seite und kauerte sich zusammen.

Da fiel sein Blick auf eine dunkle Gestalt, welche in einiger Entfernung auf einem Stein saß und ihn unter einer schweren Kapuze heraus interessiert beobachtete.

Varkur klaubte sich zusammen und stemmte seinen schmerzenden Menschenkörper in die Höhe. Die dunkle Gestalt erhob sich zeitgleich mit einer einzigen fließenden Bewegung von ihrem Stein und griff nach etwas Langem neben ihr. Die Luft flimmerte, als wäre es Hochsommer, dabei war das Graue Gebirge genauso von Eis und Schnee erfüllt, wie es Hadria beim Angriff der Mächte des Meeres gewesen war.

Varkurs Herz füllte sich zu dem Grade mit Furcht, zu dem der Dunkle Magier überhaupt noch zu einer solchen Gefühlsregung imstande war. Denn an der Seite dieser Gestalt erkannte er einen grünlich schimmernden Hammer, der eigentlich im Hohen Norden im Eisernen Turm eingeschlossen sein sollte. Er kannte den Namen dieses Hammers, und er kannte den Namen desjenigen, der ihn selbst in die Esse gehalten hatte. Orweyn, der mächtigste aller Zauberer.

Die Hadrische Unterwelt wandelte alles, was durch sie wanderte. Einst war Orweyn mit drei gewöhnlichen Waffen und einem Falken in die Hadrische Unterwelt gegangen und mit drei magischen Waffen und dem riesigen Falken von Yra zurückgekehrt. Doch an Orweyn war den Legenden nach niemandem eine Veränderung aufgefallen. Konnte es sein, dass die Unterwelt auch ihn gestärkt hatte? Dass sie ihm Ewiges Leben verliehen hatte oder sonst erlaubt hatte, die langen Jahrzehnte im Eisernen Turm zu überleben?

Damals, als der junge, schwache Varkur versucht hatte, an diesen Hammer zu kommen, damals, als er das Siegel zum Eisernen Turm gesprengt hatte, war da jemand aus dem Turm geschlüpft? Konnte es sein, dass der Dunkle Magier nun tatsächlich vor Orweyn stand?

Varkur ballte seine Echsenklauen zu Fäusten und rief im Geiste seinen Stab, doch dieser befand sich wohl immer noch unter Tonnen von Stein begraben in den Tiefen des unterirdischen Zwergennetzwerks.

„Mein Qurun!“, rief der schimmernde Orweyn theatralisch, „Das Ende Hadrias, wie ich es vorhergesehen hatte! Wenn Feuer und Turm miteinander ringen, so wirst du sie bezwingen!“

Varkur wusste, was Orweyn in seinen letzten überlieferten Lebenstagen angerichtet hatte, welcher Aufgabe er sich verschrieben hatte. Wenn dieser Magier noch lebte, so nur, um die Spuren der Dunklen Magie ein für alle Mal auszulöschen. Keine sonderlich schöne Aussicht für einen Dunklen Magier.

Noch während Varkur versuchte, sein Gleichgewicht unter Kontrolle zu kriegen und sich kampfbereit hinzustellen, die Schattenseelen wieder hervorzurufen, verschwand das seltsame Schimmern um die fremde Gestalt und enthüllte, dass sie nicht Orweyns Hammer der Stärke in ihrer Hand trug, sondern bloß einen Holzstab, wie ihn die Zauberer von Hadria ihren Speichelleckern austeilten, sobald sie sie in ihren Augen für würdig empfangen.

Und als die Gestalt ihre Kapuze zurückschlug, erkannte Varkur auch, dass das nicht Orweyn sein konnte. Die Person war glatzköpfig wie er, doch trug sie eine Augenbinde und ihre Haut schimmerte bläulich im Licht der aufgehenden Sonne. Das war Leander, der Seher von Narkon, der Hademars Seele in diesen kleinen weißen Kristall eingesperrt hatte. Dem Kristall, der nun entweder unter Tonnen von Stein begraben war oder sich in den Händen von Nomion oder dieser Schattenhexe befand. Nichts davon war eine schöne Aussicht. Und nichts davon würde es Varkur leicht machen, ihn wieder aufzuspüren. Immerhin konnte er sich nun diesen elenden Seher vorknöpfen, ohne dass ihm ein einen kolossalen Kollateralschaden in Kauf nehmender kauziger Kräuterkundler in die Sternkrautsuppe spucken würde.

„Du!“, rief Varkur aus und rief seinen dunklen Nebel hervor.

„Haltet ein, o Qurun“, sprach Leander mit einer Stimme, die nun wieder nach dem Seher und nicht nach dem heiseren Orweyn klang.

Qurun. Das Wort hallte in Varkurs Kopf nach. Er mochte diesen Klang.

„Wir müssen uns überhalten, o Qurun“, fuhr der Seher fort, „Über die elenden hadrischen Zaubererorden. Über die hirnlosen Helden von Andor. Und über Krahd.“

Varkur blieb unentschlossen stehen. Dann überbrückte er die Distanz zwischen sich und dem Seher mit einem gewaltigen Satz und sprach zum erschrockenen Seher:

„Sprich.“

\begin{center}
    Weiter geht es in \hypref{Der Giftzwerg und das Drachenherz (2023)}.
\end{center}












\newpage
\section{Epilog: Eines Hexers Abschied}



Die CAPELLA stach langsam vom Hafen am Wachsamen Wald in die Hadrische See. Meres, der Hexer von Andor, stand an der Reling und blickte gen Norden. Die See war unerwartet ruhig, fast als würde sie ihn erwarten. Nebelfetzen wogten über die Wasseroberfläche. Hier und dort konnte er vereinzelt Nixenschwänze ausmachen, als die scheuen Wasserbewohner vor dem Pfad des Schiffs auswichen .

Meres schnippte unruhig vor sich hin, woraufhin grüne Flämmchen über seine Handfläche flackerten. Er hatte sich nicht an Rekas Wunsch gehalten und brach nun in den Norden auf, ohne vorher seine einstige Lehrmeisterin aufzusuchen. Nunmehin war es schon über vierzehn Jahre her, dass Meres sich im Zorn von seiner einstigen Lehrmeisterin abgewandt hatte. Ihre verschiedenen Sichten der Welt und der Hexenkunst hatten schon einige Male zu Spannungen zwischen den beiden geführt, und die direkte Art der beiden hatte diese Spannungen kaum senken können. Sie hatten auch ihre guten Momente gehabt, wenn sie über einem neuen Trank getüftelt oder ein seltenes Kraut aus dem Barbarenland untersucht hatten. Meres wusste auch, dass die guten Momente die schlechten aufgewogen hatten. Dennoch hatten die beiden sich zu oft gegenseitig nicht ausgehalten.

Und dann war es geschehen. Als Reka verlangt hatte, er müsse eines seiner feurigeren Experimente sofort beenden, hatte Meres zu widersprechen begonnen und die Hexe kurzerhand dieses einfach selbst zu löschen versucht. Dabei konnte man diese Form der Glut nicht mit Wasser löschen. Im Gegenteil. Das hätte sie doch wissen müssen! Aber wie man in Andor so schön sagte: Selbst Brandur war nicht gegen Fehler gefeit.

Grünes Feuer war alles, an das Meres sich danach noch erinnerte. Grünes Feuer war überall gewesen und hatte die Hütte der Hexe vollends niedergebrannt. Die Kinder der Umgebung, die am nächsten Tag bei der Hütte aufkreuzten, um von Reka Lesen, Schreiben und Rechnen (von vielen Eltern abschätzig „Bewahrer-Zeug“ genannt) beigebracht zu bekommen, hatten nur noch grünlich schwelende Ascheresten vorgefunden. Und sowohl die Hexe als auch Meres hatten einander die Schuld in die Stiefel geschoben. Nicht, dass das inzwischen noch relevant wäre.

Jahrelang hatte sich Meres von Andor, dem Königreich der Hexe, ferngehalten. Hin und wieder war er zurückgekehrt, aus Faszination oder einem momentanen Wunsch. Zwischendurch hatte er sogar Reka einige Male erblickt gehabt und die Gefühle des Zorns und des Schams waren jedes Mal erneut hochgekommen, wenn auch immer schwächer. Und so hatte sich Meres immer wieder ab- und neuen Zielen zugewandt.

Neuen Zielen wie dieser mysteriösen Nebelinsel Narkon, von der der Seher Leander Meres erzählt hatte. Meres hatte gefühlt, dass der zwielichtige Seher irgendwelche Hintergedanken hatte, gar eine persönliche emotionale Bindung dazu, was mit dieser Insel geschah. Meres hatte ihm klargemacht, dass er sich nicht auf eine Reise in den Norden begeben würde, ohne zu wissen, warum er dorthin geschickt würde. Und dass er es nicht mochte, hintergangen zu werden.

Da hatte Leander ihm die Wahrheit aufgetischt und beschämt gesagt, dass ein Fluch über der Insel lag und dass eine für Leander sehr wichtige Person dort festsaß, vermutlich, ohne sich überhaupt an ihn zu erinnern. Meres hatte Mitleid gekriegt. Er verfügte über Tränke, Salben und Kräuter, von denen niemand sonst Bescheid wusste. Wer, wenn nicht er, könnte diesen Fluch zu brechen versuchen? Wer, wenn nicht er, würde sich den Gefangenen des Fluches annehmen?

„Meres! So warte doch!“

Meres erstarrte und wurde aus seinen Gedanken gerissen. Das war Rekas Stimme gewesen, und sie kam von hinten, vom andorischen Ufer. Die raue Stimme, die er so oft als zeternd in Erinnerung hatte, klang nun wieder weich und einladend. Er haderte mit sich selbst. Die CAPELLA stach gen Norden, bald schon würde sie zu weit entfernt von der Küste sein, als dass er sich noch sinnvoll mit Reka unterhalten könnte. Er könnte sie wieder vergessen, das alles hinter sich lassen. Für seine Gemütslage wäre das wahrscheinlich am besten.

Aber er könnte auch die Gelegenheit nutzen, sich der Hexe zu stellen.

Meres machte abrupt kehrt, rannte an der Seite des Schiffes entlang zum Heck und spähte in Richtung des Wachsamen Waldes. Im grauen Licht des ersten Morgens wirkten die letzten Baumstämme des Waldrandes schwarz wie die Stäbe eines Käfigs. Und vor diesem Käfig stand sie.

Reka.

Sie trug ihren üblichen zerrissenen grauen Mantel und winkte hektisch dem abziehenden Schiff hinterher. Der Wind trug ihre krächzende Stimme klar und deutlich an Meres‘ Ohren.

„Meres! Halte ein! Es tut mir leid! Ich will dich nicht vertreiben!“

Sie hielt inne. Meres stockte ebenfalls. Selbst die CAPELLA selbst blieb plötzlich stehen. Meres blickte überrascht zu Kapitänin Mondrianne hinüber, welche neben ihrer Tochter am Steuer saß.

Mondrianne meinte belustigt: „Tu dir keinen Zwang an, Hexer, aber ich hab die alte Reka meiner Lebtag noch nie so schnell anrennen sehen. Ich will wissen, wie das ausgeht. Die CAPELLA fährt erst weiter, wenn ihr eure Angelegenheit geregelt habt.“

Meres‘ Blick wanderte wieder zu Reka. Er wusste nicht, was er antworten sollte. Er empfand schon lange keine starken Emotionen gegenüber Reka mehr, oder zumindest redete er sich das ein. Auch ihm tat es leid, was zwischen ihnen vorgefallen war.

„Ich habe einen Fehler begangen“, schrie Reka weiter. Kurz hielt sie inne, wohl ob der Mehrdeutigkeit ihres Zugeständnisses, dann meinte sie: „Ich habe dem Seher Leander von dir erzählt. Er manipuliert dich. Narkon ist eine Falle!“

„Bin ich zu weit weg, um meine Gedanken zu lesen?“, rief Meres ruhig zurück, „Meine Entscheidung steht fest.“

„Ach ja. Deine Meinung hat sich schon immer verhärtet, wenn man dir mit Rat zu kommen versucht. Aber du wirst es alleine schlicht nicht schaffen, dieser Falle zu entkommen.“

„Immer meinst du, die Zukunft bereits zu kennen, Reka. Keiner kann das mit Sicherheit. Ich will zumindest versuchen, diesen leidenden Wesen zu helfen. Das ist das, was ein Held tun würde.“

Reka blieb still.

Warum hatte die Hexe ihr Wiedersehen so publik machen müssen? Andererseits gaben ihm die hinüberspähenden Zuschauer auf dem Schiff zumindest einen guten Grund, das Gespräch schnell hinter sich zu bringen. Und was hatten Reka und er sich auch groß zu sagen? Sie wussten beide, was damals vorgefallen war. Sie kamen nicht groß miteinander klar, aber sie hassten einander auch nicht mehr. Dafür war zu viel Wasser die Narne hinuntergegangen in der Zwischenzeit. Jetzt war seine Gelegenheit, einen Schlussstrich zu setzen.

„Auch mir tut es leid“, sprach Meres schlicht. War das eine gute Antwort? Eine sinnvolle? Eine hilfreiche?

Stille. Er wandte seinen Blick ab und schnippte wieder unruhig vor sich hin, woraufhin kleine kalte grüne Flämmlein an seinen Fingerspitzen auftanzten.

Rekas Stimme ertönte, und zum ersten Mal glaubte Meres, wieder ein bisschen von der Zuneigung zu spüren, die Reka ihm früher immer entgegengebracht hatte.

„Mir bleibt wohl nichts anderes übrig, als dir viel Glück zu wünschen auf deinem weiteren Weg“, sprach Reka sanft, „Das ist eigentlich auch das, weswegen ich dich überhaupt erst wiedersehen wollte. Ich will...“

Die sonst so selbstsichere Reka stockte an den Worten, als sie fortfuhr: „Ich will nur ausdrücken, dass du Andor nicht mehr meinetwegen meiden musst. Wir haben beide Fehler begangen, aber diese sind lange her. Wir können sie hinter uns lassen, wenn wir das beide wollen. Ich werde hier sein, wenn du Lust hast, von deinen Reisen und deinen Erkenntnissen zu berichten. Es würde mich sogar freuen Das nächste Mal, wenn du zurückkommst. Falls. Du verstehst schon. Viel Glück mit Narkon, Meres. Du wirst es brauchen.“

Meres wich weiterhin ihrem Blick aus.

Doch er lächelte.

